\section{AMC 12 12B 2022}
        \begin{enumerate}
        	\item (\textit{AMC 12 12B 2022}) Define $x\diamond y$ to be $|x-y|$ for all real numbers $x$ and $y.$ What is the value of 
 \[(1\diamond(2\diamond3))-((1\diamond2)\diamond3)?\]
$\textbf{(A)}\ {-}2 \qquad \textbf{(B)}\ {-}1 \qquad \textbf{(C)}\ 0 \qquad \textbf{(D)}\ 1 \qquad \textbf{(E)}\ 2$

 \textbf{Answer: }A



\textbf{Solution 1}

We have \begin{align}(1\diamond(2\diamond3))-((1\diamond2)\diamond3) &= |1-|2-3|| - ||1-2|-3| \\&= |1-1| - |1-3| \\&= 0-2 \\&= \boxed{\textbf{(A)}\ {-}2}.\end{align}~MRENTHUSIASM

 



\textbf{Solution 2}

Observe that the $\diamond$ function is simply the positive difference between two numbers. Thus, we evaluate: the difference between $2$ and $3$ is $1;$ the difference between $1$ and $1$ is $0;$ the difference between $1$ and $2$ is $1;$ the difference between $1$ and $3$ is $2;$ and finally, $0-2=\boxed{\textbf{(A)}\ {-}2}.$



	\item (\textit{AMC 12 12B 2022}) In rhombus $ABCD$, point $P$ lies on segment $\overline{AD}$ so that $\overline{BP}$ $\perp$ $\overline{AD}$, $AP = 3$, and $PD = 2$. What is the area of $ABCD$? (Note: The figure is not drawn to scale.)


 
                    \begin{figure}[H]
                        \centering
                        \includegraphics[width=0.5\textwidth]{Problem/AMC_12_12B_2022_2_Problem_Figure_1.png}
                    \end{figure}
                    
$\textbf{(A) }3\sqrt 5 \qquad \textbf{(B) }10 \qquad \textbf{(C) }6\sqrt 5 \qquad \textbf{(D) }20\qquad \textbf{(E) }25$

 \textbf{Answer: }D



\textbf{Solution}

Notice that $BC = AD = AP + PD = 3 + 2 = 5$.Now, because the question tells us $ABCD$ is a rhombus, $BC = AB = 5$. Thus, by Pythagorean Theorem on $APB$, $BP = 4$.Therefore, 
 \[[ABCD] = AD \times BP = 5 \times 4 = \boxed{\text{D 20}}\]~DUODUOLUO(First time using $\LaTeX$, may have some minor mistakes)

 


	\item (\textit{AMC 12 12B 2022}) How many of the first ten numbers of the sequence $121, 11211, 1112111, \ldots$ are prime numbers?


 $\textbf{(A) } 0 \qquad \textbf{(B) }1 \qquad \textbf{(C) }2 \qquad \textbf{(D) }3 \qquad \textbf{(E) }4$

 \textbf{Answer: }A



\textbf{Solution 1 (Generalization)}

The $n$th term of this sequence is
 \[\sum_{k=n}^{2n}10^k + \sum_{k=0}^{n}10^k = 10^n\sum_{k=0}^{n}10^k + \sum_{k=0}^{n}10^k = \left(10^n+1\right)\sum_{k=0}^{n}10^k.\]It follows that the terms are
 \begin{align*} 121 &= 11\cdot11, \\ 11211 &= 101\cdot111, \\ 1112111 &= 1001\cdot1111, \\ & \ \vdots \end{align*}
Therefore, there are $\boxed{\textbf{(A) } 0}$ prime numbers in this sequence.

 



\textbf{Solution 2 (Detailed Explanation of Solution 1)}

Denote this sequence as $a_{n}$, then we can find that
 \begin{align*} a_{1} &= 121 = 10^2 + 2\cdot10 + 1 = (10^2 + 10) + (10 + 1), \\ a_{2} &= 11211 = (10^4 + 10^3 + 10^2) + (10^2 + 10 + 1), \\ a_{3} &= 1112111 = (10^6 + 10^5 + 10^4 + 10^3) + (10^3 + 10^2 + 10 + 1), \\ & \ \vdots \end{align*}
So, we can induct that the general term is
 \begin{align*} a_n &= (10^{2n} + 10^{2n-1} + \ldots + 10^{n+1} + 10^n) + (10^n + 10^{n-1} + \ldots +10 + 1) \\ &= 10^n\cdot(10^n + 10^{n-1} + \ldots +10 + 1) + (10^n + 10^{n-1} + \ldots +10 + 1) \\ &= \left(10^n+1\right)\sum_{k=0}^{n}10^k. \end{align*}
Therefore, there are $\boxed{\textbf{(A) } 0}$ prime numbers in this sequence.

 



\textbf{Solution 3 (Simple Sums)}

Observe how 
 \begin{align*} 121 &= 110 + 11, \\ 11211 &= 11100 + 111, \\ 1112111 &= 1111000 + 1111, \\ & \ \vdots \end{align*}
all take the form of 
 \[\underbrace{111\ldots}_{n+1}\underbrace{00\ldots}_{n} + \underbrace{111\ldots}_{n+1} = \underbrace{111\ldots}_{n+1}(10^{n} + 1).\]Factoring each of the sums, we have 
 \[11(10+1), 111(100+1), 1111(1000+1), \ldots\] respectively. With each number factored, there are $\boxed{\textbf{(A) } 0}$ primes in the set.

 



\textbf{Solution 4 (Educated Guess)}

Note that $121$ is divisible by $11$ and $11211$ is divisible by $3$. Because this is Problem 6 of the AMC 10, we assume we do not need to check two-digit prime divisibility or use obscure theorems. Therefore, the answer is $\boxed{\textbf{(A) } 0}.$



	\item (\textit{AMC 12 12B 2022}) For how many values of the constant $k$ will the polynomial $x^{2}+kx+36$ have two distinct integer roots?


 $\textbf{(A) }6 \qquad \textbf{(B) }8 \qquad \textbf{(C) }9 \qquad \textbf{(D) }14 \qquad \textbf{(E) }16$

 \textbf{Answer: }B



\textbf{Solution 1}

Let $p$ and $q$ be the roots of $x^{2}+kx+36.$ By \href{https://artofproblemsolving.com/wiki/index.php?title=Vieta\%27s\_Formulas}{Vieta's Formulas}, we have $p+q=-k$ and $pq=36.$

 It follows that $p$ and $q$ must be distinct factors of $36.$ The possibilities of $\{p,q\}$ are 
 \[\pm\{1,36\},\pm\{2,18\},\pm\{3,12\},\pm\{4,9\}.\]Each unordered pair gives a unique value of $k.$ Therefore, there are $\boxed{\textbf{(B) }8}$ values of $k,$ corresponding to $\mp37,\mp20,\mp15,\mp13,$ respectively.

 



\textbf{Solution 2}

Note that $k$ must be an integer. Using the \href{https://artofproblemsolving.com/wiki/index.php?title=Quadratic\_formula}{quadratic formula}, $x=\frac{-k \pm \sqrt{k^2-144}}{2}.$ Since $4$ divides $144$ evenly, $k$ and $k^2-144$ have the same parity, so $x$ is an integer if and only if $k^2-144$ is a perfect square.

 Let $k^2-144=n^2.$ Then, $(k+n)(k-n)=144.$ Since $k$ is an integer and $144$ is even, $k+n$ and $k-n$ must both be even. Assuming that $k$ is positive, we get $5$ possible values of $k+n$, namely $2, 4, 8, 6, 12$, which will give distinct positive values of $k$, but $k+n=12$ gives $k+n=k-n$ and $n=0$, giving $2$ identical integer roots. Therefore, there are $4$ distinct positive values of $k.$ Multiplying that by $2$ to take the negative values into account, we get $4\cdot2=\boxed{\textbf{(B) }8}$ values of $k$.

 



\textbf{Solution 3 (Pythagorean Triples)}

Proceed similar to Solution 2 and deduce that the discriminant of $x^{2}+kx+36$ must be a perfect square greater than $0$ to satisfy all given conditions. Seeing something like $k^2-144$ might remind us of a right triangle, where $k$ is the hypotenuse, and $12$ is a leg. There are four ways we could have this: a $9$-$12$-$15$ triangle, a $12$-$16$-$20$ triangle, a $5$-$12$-$13$ triangle, and a $12$-$35$-$37$ triangle.

 Multiply by $2$ to account for negative $k$ values (since $k$ is being squared), and our answer is $\boxed{\textbf{(B) }8}$.

 



\textbf{Solution 4}

Since $36 = 2^2\cdot3^2$, that means there are $(2+1)(2+1) = 9$ possible factors of $36$. Since $6 \cdot 6$ violates the distinct root condition, subtract $1$ from $9$ to get $8$. Each sum is counted twice, and we count of those twice for negatives. This cancels out, so we get $\boxed{\textbf{(B) }8}$.

 



\textbf{Solution 5 (Discriminant and Factor Rainbow)}

Since to have two solutions, the discriminant must be greater than $0$, $k^2 - 144 > 0$ so k must be greater than $12$ or less than $-12$(since it's squared). Next, we know that when factoring, the two numbers must multiply to $36$. So you can create a factor rainbow(remember from elementary school) which has the pairs of integers that multiply to $36( 1, 2, 3, 4, 6, 6, 9, 12, 18, 36 )$ (you just look at the position and the one on the opposite side, for example $2$ is the second value, and $18$ is the second last value, and they multiply to 36). See that there are $4$ pairs of values which add up to greater than $12$($6+6$ is equal to $12,$ not greater), and we look at the negative values of them too(since k is squared) we get $\boxed{\textbf{(B) }8}$ total values.

 


	\item (\textit{AMC 12 12B 2022}) The point $(-1, -2)$ is rotated $270^{\circ}$ counterclockwise about the point $(3, 1)$. What are the coordinates of its new position?


 $\textbf{(A) }\ (-3, -4) \qquad \textbf{(B) }\ (0,5) \qquad \textbf{(C) }\ (2,-1) \qquad \textbf{(D) }\ (4,3) \qquad \textbf{(E) }\ (6,-3)$

 \textbf{Answer: }B



\textbf{Solution 1 (Cartesian Plane)}

$(-1,-2)$ is $4$ units west and $3$ units south of $(3,1)$. Performing a counterclockwise rotation of $270^{\circ}$, which is equivalent to a clockwise rotation of $90^{\circ}$, the answer is $3$ units west and $4$ units north of $(3,1)$, or $\boxed{\textbf{(B)}\ (0,5)}$.

 



\textbf{Solution 2 (Complex Numbers)}

We write $(-1, -2)$ as $-1-2i.$ We'd like to rotate about $(3, 1),$ which is $3+i$ in the complex plane, by an angle of $270^{\circ} = \frac{3\pi}{2} \text{ rad}$ counterclockwise.

 The formula for rotating the complex number $z$ about the complex number $w$ by an angle of $\theta$ counterclockwise is given as $(z-w)e^{\theta i} + w.$ Plugging in our values $z = -1-2i, w = 3+i, \theta = \frac{3\pi}{2}$, we evaluate the expression as $((-1-2i) - (3+i))e^{\frac{3\pi i}{2}} + (3+i) = (-4-3i)(-i) + (3+i) = 4i-3i^2+3+i = 4i-3+3+i = 5i,$ which corresponds to $\boxed{\textbf{(B)}\ (0,5)}$ on the Cartesian plane.

 


 



\textbf{Solution 3 (Cheese)}

A $270^{\circ}$ rotation counterclockwise is equivalent to a $90^{\circ}$ rotation clockwise, which means the resultant point should be around quadrant 2. Looking at the answer choices the only option that works in $\boxed{\textbf{(B)}\ (0,5)}$.

 


	\item (\textit{AMC 12 12B 2022}) Consider the following $100$ sets of $10$ elements each:

 \begin{align*} &\{1,2,3,\ldots,10\}, \\ &\{11,12,13,\ldots,20\},\\ &\{21,22,23,\ldots,30\},\\ &\vdots\\ &\{991,992,993,\ldots,1000\}. \end{align*}

How many of these sets contain exactly two multiples of $7$?


 $\textbf{(A)}\ 40\qquad\textbf{(B)}\ 42\qquad\textbf{(C)}\ 43\qquad\textbf{(D)}\ 49\qquad\textbf{(E)}\ 50$

 \textbf{Answer: }B



\textbf{Solution 1 (Competition)}

There are $\left\lfloor \frac{1000}{7} \right\rfloor = 142$ numbers divisible by 7. We split these into 100 sets containing 10 numbers each, giving us 1.42 multiples of 7 per set. After the first set, the numbers come evenly, and we multiply 100 by $1.42 - 1 = \boxed{\textbf{(B)}\ 42}.$

 (Genuinely genius solution ~ MW)

 



\textbf{Solution 2 (Casework)}

We apply casework to this problem. The only sets that contain two multiples of seven are those for which:

 
            \begin{enumerate} 
        \item The multiples of $7$ are $1\pmod{10}$ and $8\pmod{10}.$ That is, the first and eighth elements of such sets are multiples of $7.$The first element is $1+10k$ for some integer $0\leq k\leq99.$ It is a multiple of $7$ when $k=2,9,16,\ldots,93.$

        \item The multiples of $7$ are $2\pmod{10}$ and $9\pmod{10}.$ That is, the second and ninth elements of such sets are multiples of $7.$The second element is $2+10k$ for some integer $0\leq k\leq99.$ It is a multiple of $7$ when $k=4,11,18,\ldots,95.$

        \item The multiples of $7$ are $3\pmod{10}$ and $0\pmod{10}.$ That is, the third and tenth elements of such sets are multiples of $7.$The third element is $3+10k$ for some integer $0\leq k\leq99.$ It is a multiple of $7$ when $k=6,13,20,\ldots,97.$

            \end{enumerate}
            Each case has $\left\lfloor\frac{100}{7}\right\rfloor=14$ sets. Therefore, the answer is $14\cdot3=\boxed{\textbf{(B)}\ 42}.$




\textbf{Solution 3 (Find A Pattern)}

We find a pattern. 
 \begin{align*} &\{1,2,3,\ldots,10\}, \\ &\{11,12,13,\ldots,20\},\\ &\{21,22,23,\ldots,30\},\\ &\vdots\\ &\{991,992,993,\ldots,1000\}. \end{align*}
We can figure out that the first set has $1$ multiple of $7$. The second set also has $1$ multiple of $7$. The third set has $2$ multiples of $7$. The fourth set has $1$ multiple of $7$. The fifth set has $2$ multiples of $7$. The sixth set has $1$ multiple of $7$. The seventh set has $1$ multiple of $7$. The eighth set has $2$ multiples of $7$. Disregarding the first set and then calculating this pattern further, we can see (reasonably) that it repeats for each $1$ sets.We see that the pattern for the number of multiples per $7$ sets (again, disregarding the first set) goes: $1,2,1,2,1,1,2.$ So, for every $7$ sets after the first, there are three sets with $2$ multiples of $7$. We calculate $\left\lfloor\frac{100-1}{7}\right\rfloor$ and multiply that by $3$. (We also disregard the remainder of $1$ since it doesn't add any extra sets with $2$ multiples of $7$.). We get $14\cdot3= \boxed{\textbf{(B) }42}$.

 



\textbf{Solution 4 (Fastest)}

Each set contains exactly $1$ or $2$ multiples of $7$ (Because it will need at least $2$ * $7$ = $14$ numbers in a group for $3$ or more multiples of $7$).

 There are $\dfrac{1000}{10}=100$ total sets and $\left\lfloor\dfrac{1000}{7}\right\rfloor = 142$ multiples of $7$. Thus, there are $142-100=\boxed{\textbf{(B) }42}$ sets with $2$ multiples of $7$.

 



\textbf{Solution 5 (Simple Counting, Similar to Solution 1)}

Consecutive multiples of $7$ must differ by $7$. So, if a set $\{\ldots1,\ldots2,\ldots3,(\ldots),\ldots8,\ldots9,\ldots0\}$ contains two multiples of $7$, they must end with the digits $1$ and $8$, $2$ and $9$, or $3$ and $0$. This reduces the problem to counting the amount of multiples of $7$ less than $1000$ that end with $1$, $2$, and $3$.

 The first multiple of $7$ that ends with $1$ is $21$. The next multiple that ends with $1$ occurs $70$ later, since that is the smallest multiple of $7$ we can add to $21$ without affecting the last digit. The greatest number of $70$'s we can add to $21$ while keeping it less than $1000$ is $13$, because $21 + 70(13) = 931$. Therefore, the set of multiples of $7$ less than $1000$ ending with $1$ is $\{21 + 70(0), 21+70(1),\ldots,21+70(13)\}$, meaning there are $14$ of these particular multiples. We can use the same reasoning to count the multiples of $7$ that end with $2$ and $3$. 

 The first multiple of $7$ that ends with $2$ is $42$. The greatest number of $70$'s we can add to $42$ here is also $13$, since $42 + 70(13) = 952$. The set of multiples of $7$ less than $1000$ ending with $2$ is $\{42 + 70(0), 42+70(1),\ldots,42+70(13)\}$, giving $14$ multiples.

 The first multiple of $7$ that ends with $3$ is $63$. The greatest number of $70$'s we can add to $63$ here is yet again $13$, since $63 + 70(13) = 973$. The set of multiples of $7$ less than $1000$ ending with $3$ is $\{63 + 70(0), 63+70(1),\ldots,63+70(13)\}$, giving another $14$ multiples.

 In total, there are $14 + 14 + 14 = 42$ of these multiples, and so $\boxed{\textbf{(B) }42}$ sets with two multiples of $7$.

 



\textbf{Solution 6 (Very Fast System of Equations)}

Let $a$ be the number of sets with $1$ multiple of $7$ and $b$ be the number of sets with $2$ multiples of $7$. Note that it is impossible for a set to have more than two multiples of $7$. 

 Since there are a total of $100$ sets, $a+b=100$. Also, since there are $\lfloor \frac{1000}{7}\rfloor = 142$ multiples of $7$ between $1$ and $1000$, we must have $a+2b=142$. 

 Solving the system of equations

 
 \[a+b=100\]
 \[a+2b=142\]

 for $b$ gives us the answer of $\boxed{\textbf{(B) }42}$.

 



\textbf{Solution 7 (reasonable count)}

(Similar to Solution 3, but a little more intuitive and less numbers.) Note that this system loops every cycle of length $70$, or $7$ such sets. From $1$ to $70$, the multiples of $7$ are $7$, $14$, $21$, $28$, $35$, $42$, $49$, $56$, $63$, and $70$; note that $(21,28),(42,49),$ and $(63,70)$ are in the same sets. Thus, for every $7$ sets, we have $3$ sets with exactly two multiples of $7$. We have $100$ sets, which is $98+2=7\cdot14+2$; the first $98$ sets contain $\dfrac37(7\cdot14)=42$ desired sets. The last two sets comprise the integers from $981$ to $999$; the multiples of $7$ here are $987$ and $994$. Neither of the two last sets contains two multiples of $7$, so our answer is simply $\boxed{\textbf{(D) }42}$. 

 



\textbf{Solution 8 (Pigeonhole Principle, very fast)}

Consider any consecutive seven sets beginning on one more than a multiple of $70$. There are guaranteed to be $10$ multiples of $7$ in these seven sets. Note that no set may have three multiples of $7$. This means that some three of the ten sets have $2$ multiples of seven. There are $14$ such groups of seven sets, and then the sets $\{981, 982, \dots, 989\}$ and $\{991, 992, \dots, \}$. These last two sets are congruent to $\{1, 2, \dots, 9\}$ and $\{11, 12, \dots, 19\}$ modulo 7, thus making it easy to check that each only has one multiple of $7$. Therefore, there are $14 \cdot 3 = \boxed{\textbf{(D) }42}$ sets with $2$ multiples of $7$. 

 


	\item (\textit{AMC 12 12B 2022}) Camila writes down five positive integers. The unique mode of these integers is $2$ greater than their median, and the median is $2$ greater than their arithmetic mean. What is the least possible value for the mode?


 $\textbf{(A) }5\qquad\textbf{(B) }7\qquad\textbf{(C) }9\qquad\textbf{(D) }11\qquad\textbf{(E) }13$

 \textbf{Answer: }D



\textbf{Solution 1 (Variables)}

Let $M$ be the median. It follows that the two largest integers are both $M+2.$

 Let $a$ and $b$ be the two smallest integers such that $a<b.$ The sorted list is 
 \[a,b,M,M+2,M+2.\]Since the median is $2$ greater than their arithmetic mean, we have $\frac{a+b+M+(M+2)+(M+2)}{5}+2=M,$ or 
 \[a+b+14=2M.\]Note that $a+b$ must be even. We minimize this sum so that the arithmetic mean, the median, and the unique mode are minimized. Let $a=1$ and $b=3,$ from which $M=9$ and $M+2=\boxed{\textbf{(D)}\ 11}.$




\textbf{Solution 2 (Elimination)}

We can also easily test all the answer choices. (This strategy is generally good to use for multiple-choice questions if you don't have a concrete method to proceed with!)

 For answer choice $\textbf{(A)},$ the mode is $5,$ the median is $3,$ and the arithmetic mean is $1.$ However, we can quickly see this doesn't work, as there are five integers, and they can't have an arithmetic mean of $1$ while having a mode of $5.$ 

 Trying answer choice $\textbf{(B)},$ the mode is $7,$ the median is $5,$ and the arithmetic mean is $3.$ From the arithmetic mean, we know that all the numbers have to sum to $15.$ We know three of the numbers: $\underline{\hspace{3mm}},\underline{\hspace{3mm}},5,7,7.$ This exceeds the sum of $15.$

 Now we try answer choice $\textbf{(C)}.$ The mode is $9,$ the median is $7,$ and the arithmetic mean is $5.$ From the arithmetic mean, we know that the list sums to $25.$ Three of the numbers are $\underline{\hspace{3mm}},\underline{\hspace{3mm}},7,9,9,$ which is exactly $25.$ However, our list needs positive integers, so this won't work.

 Since we were really close on answer choice $\textbf{(C)},$ we can intuitively feel that the answer is probably going to be $\textbf{(D)}.$ We can confirm this by creating a list that satisfies the problem and choose $\textbf{(D)}: 1,3,9,11,11.$

 So, our answer is $\boxed{\textbf{(D)}\ 11}.$



	\item (\textit{AMC 12 12B 2022}) What is the graph of $y^4+1=x^4+2y^2$ in the coordinate plane?


 $\textbf{(A) }\ \text{two intersecting parabolas} \qquad \textbf{(B) }\ \text{two nonintersecting parabolas} \qquad \textbf{(C) }\ \text{two intersecting circles} \qquad$
$\textbf{(D) }\ \text{a circle and a hyperbola} \qquad \textbf{(E) }\ \text{a circle and two parabolas}$

 \textbf{Answer: }D



\textbf{Solution 1}

Since the equation has even powers of $x$ and $y$, let $y'=y^2$ and $x' = x^2$. Then $y'^2 + 1 = x'^2 + 2y'$. Rearranging gives $y'^2 - 2y' + 1 = x'^2$, or $(y'-1)^2=x'^2$. There are two cases: $y' \leq 1$ or $y' > 1$.

 If $y' \leq 1$, taking the square root of both sides gives $1 - y' = x'$, and rearranging gives $x' + y' = 1$. Substituting back in $x'=x^2$ and $y'=y^2$ gives us $x^2+y^2=1$, the equation for a circle.

 Similarly, if $y' > 1$, we take the square root of both sides to get $y' - 1 = x'$, or $y' - x' = 1$, which is equivalent to $y^2 - x^2 = 1$, a hyperbola.

 Hence, our answer is $\boxed{\textbf{(D)}\ \text{a circle and a hyperbola}}$.

 (Solutions 1 and 2 are in essence the same; Solution 1 lets $(x',y')=\left(x^2,y^2\right)$ for convenience, but the two solutions are otherwise identical.) 

 



\textbf{Solution 2 (Factoring)}

We can subtract $2y^2$ from both sides and factor the left side to get $(y^2-1)^2 = x^4$. If we take the square root of both sides, we are left with two equations: $y^2-1 = x^2$ and $y^2-1 = -x^2$. In the former case, we get $y^2-x^2=1$, which, is the formula for a hyperbola. This means that a hyperbola will appear in the graph of the equation. In the latter case, we get $x^2+y^2=1$, which is the equation of a circle which will also appear in the graph. Looking at our options, this is the only valid answer. Hence, the answer is $\boxed{\textbf{(D)}\ \text{a circle and a hyperbola}}$.

 


	\item (\textit{AMC 12 12B 2022}) The sequence $a_0,a_1,a_2,\cdots$ is a strictly increasing arithmetic sequence of positive integers such that 
 \[2^{a_7}=2^{27} \cdot a_7.\] What is the minimum possible value of $a_2$?


 $\textbf{(A) }\ 8 \qquad \textbf{(B) }\ 12 \qquad \textbf{(C) }\ 16 \qquad \textbf{(D) }\ 17 \qquad \textbf{(E) }\ 22$

 \textbf{Answer: }B



\textbf{Solution 1}

We can rewrite the given equation as $2^{a_7-27}=a_7$. Hence, $a_7$ must be a power of $2$ and larger than $27$. The first power of 2 that is larger than $27$, namely $32$, does satisfy the equation: $2^{32 - 27} = 2^5 = 32$. In fact, this is the only solution; $2^{a_7-27}$ is exponential whereas $a_7$ is linear, so their graphs will not intersect again.

 Now, let the common difference in the sequence be $d$. Hence, $a_0 = 32 - 7d$ and $a_2 = 32 - 5d$. To minimize $a_2$, we maxmimize $d$. Since the sequence contains only positive integers, $32 - 7d > 0$ and hence $d \leq 4$. When $d = 4$, $a_2 = \boxed{\textbf{(B)}\ 12}$.

 


 


	\item (\textit{AMC 12 12B 2022}) Regular hexagon $ABCDEF$ has side length $2$. Let $G$ be the midpoint of $\overline{AB}$, and let $H$ be the midpoint of $\overline{DE}$. What is the perimeter of $GCHF$?


 $\textbf{(A) }\ 4\sqrt3 \qquad \textbf{(B) }\ 8 \qquad \textbf{(C) }\ 4\sqrt5 \qquad \textbf{(D) }\ 4\sqrt7 \qquad \textbf{(E) }\ 12$

 \textbf{Answer: }D



\textbf{Solution 1}

Let the center of the hexagon be $O$. $\triangle AOB$, $\triangle BOC$, $\triangle COD$, $\triangle DOE$, $\triangle EOF$, and $\triangle FOA$ are all equilateral triangles with side length $2$. Thus, $CO = 2$, and $GO = \sqrt{AO^2 - AG^2} = \sqrt{3}$. By symmetry, $\angle COG = 90^{\circ}$. Thus, by the Pythagorean theorem, $CG = \sqrt{2^2 + \sqrt{3}^2} = \sqrt{7}$. Because $CO = OF$ and $GO = OH$, $CG = HC = FH = GF = \sqrt{7}$. Thus, the solution to our problem is $\sqrt{7} + \sqrt{7} + \sqrt{7} + \sqrt{7} = \boxed{\textbf{(D)}\ 4\sqrt7}$.

 



\textbf{Solution 2}

Consider triangle $AFG$. Note that $AF = 2$, $AG = 1$, and $\angle GAF = 120 ^{\circ}$ because it is an interior angle of a regular hexagon. (See note for details.) 

 By the \href{https://artofproblemsolving.com/wiki/index.php?title=Law\_of\_Cosines}{Law of Cosines}, we have: 
 \begin{align*} FG^2 &= AG^2 + AF^2 - 2 \cdot AG \cdot AF \cdot \cos \angle GAF \\ FG^2 &= 1^2 + 2^2 - 2 \cdot 1 \cdot 2 \cdot \cos 120 ^{\circ} \\ FG^2 &= 5 + 4 \cdot \left( \frac 12 \right) \\ FG^2 &= 7 \\ FG &= \sqrt 7. \end{align*}
By \href{https://artofproblemsolving.com/wiki/index.php?title=SAS\_Congruence&action=edit&redlink=1}{SAS Congruence}, triangles $AFG$, $BCG$, $CDH$, and $EFH$ are congruent, and by \href{https://artofproblemsolving.com/wiki/index.php?title=CPCTC}{CPCTC}, quadrilateral $GCHF$ is a rhombus. Therefore, the perimeter of $GCHF$ is $4 \cdot FG = \boxed{\textbf{(D)}\ 4\sqrt7}$.

 Note: The sum of the interior angles of any polygon with $n$ sides is given by $180 ^{\circ} (n - 2)$. Therefore, the sum of the interior angles of a hexagon is $720 ^{\circ}$, and each interior angle of a regular hexagon measures $\frac{720 ^{\circ}}{6} = 120 ^{\circ}$.

 



\textbf{Solution 3}

We use a coordinates approach. Letting the origin be the center of the hexagon, we can let $A = (-1, \sqrt{3}), B = (1, \sqrt{3}), C = (2, 0), D = (1, -\sqrt{3}), E = (-1, -\sqrt{3}), F = (-2, 0).$ Then, $G = (0, \sqrt{3})$ and $H = (0, -\sqrt{3}).$

 We use the distance formula four times to get $CH, HF, FG, \text{ and } GC.$
 \begin{alignat*}{8} CH\textasciicircum{}2 &= (2-0)\textasciicircum{}2 + (0-(-\sqrt{3}))\textasciicircum{}2 &&= 7 &&\rightarrow CH &&= \sqrt{7}, \\ HF\textasciicircum{}2 &= (0-(-2))\textasciicircum{}2 + (-\sqrt{3}-0)\textasciicircum{}2 &&= 7 &&\rightarrow HF &&= \sqrt{7}, \\ FG\textasciicircum{}2 &= (-2-0)\textasciicircum{}2 + (0-\sqrt{3})\textasciicircum{}2 &&= 7 &&\rightarrow FG &&= \sqrt{7}, \\ GC\textasciicircum{}2 &= (0-2)\textasciicircum{}2 + (\sqrt{3}-0)\textasciicircum{}2 &&= 7 &&\rightarrow GC &&= \sqrt{7}. \end{alignat*}Thus, the perimeter of $GCHF = CH + HF + FG + GC = \sqrt{7} + \sqrt{7} + \sqrt{7} + \sqrt{7} = \boxed{\textbf{(D)}\ 4\sqrt7}$.

 Note: the last part of this solution could have been simplified by noting that $CH = HF = FG = GC = \sqrt{7}.$




\textbf{Solution 4}

Note that triangles $\triangle{GAF}, \triangle{FEH}, \triangle{HDC},$ and $\triangle{CBG}$ are all congruent, since they have side lengths of $1$ and $2$ and an included angle of $120^{\circ}.$ 

 By the Law of Cosines, 
 \[FG=\sqrt{1^2+2^2-2\cdot{1}\cdot{-\frac{1}{2}}}=\sqrt{7}.\] Therefore, $FG+GC+CH+HF=4\cdot{FG}=\boxed{\textbf{(D)}\ 4\sqrt7}.$

 -Benedict T (countmath1)

 



\textbf{Solution 5 (Answer Choices + Pythagorean Theorem Extension)}

Like the previous solutions, note that $\triangle{GAF}, \triangle{FEH}, \triangle{HDC},$ and $\triangle{CBG}$ are all congruent by SAS. It follows that quadrilateral $GCHF$ is a rhombus. 

 Recall the Pythagorean Theorem, which states $a^2+b^2=c^2$ for all right triangles, where $c$ is the hypotenuse of the triangle. However, by drawing a quick diagram of an obtuse triangle, we can see that $a^2+b^2<c^2$, in any given obtuse triangle. 

 Since $ABCDEF$ is a regular hexagon, all of its angles are obtuse. It follows that $\triangle{GAF}$ is an obtuse triangle. Using the extended Pythagorean Theorem for obtuse triangles, we have: 
 \[AG^2+AF^2<GF^2 \implies 1^2+2^2<GF^2 \implies \sqrt{5}<GF\] 

 Since $GCHF$ is a rhombus, the perimeter is $4GF$. This eliminates all answer choices but $D$ and $E$, since in all of those options $GF<\sqrt{5}$. Lastly, $E$ is eliminated due to the triangle inequality, as $1+2$ is not greater than $12/4=3$. 

 Hence, the answer is $\boxed{\textbf{(D)}\ 4\sqrt7}$.

 


	\item (\textit{AMC 12 12B 2022}) Let $f(n) = \left( \frac{-1+i\sqrt{3}}{2} \right)^n + \left( \frac{-1-i\sqrt{3}}{2} \right)^n$, where $i = \sqrt{-1}$. What is $f(2022)$?


 $\textbf{(A) }\ -2 \qquad \textbf{(B) }\ -1 \qquad \textbf{(C) }\ 0 \qquad \textbf{(D) }\ \sqrt{3} \qquad \textbf{(E) }\ 2$

 \textbf{Answer: }E



\textbf{Solution 1}

Converting both summands to exponential form, 
 \begin{align*} -1 + i\sqrt{3} &= 2e^{\frac{2\pi i}{3}}, \\ -1 - i\sqrt{3} &= 2e^{-\frac{2\pi i}{3}} = 2e^{\frac{4\pi i}{3}}. \end{align*}
Notice that the two terms in the problem are two of the third roots of unity (that is, both of them equal $1$ when raised to the power of $3$). When we replace the summands with their exponential form, we get 
 \[f(n) = \left(e^{\frac{2\pi i}{3}}\right)^n + \left(e^{\frac{4\pi i}{3}}\right)^n.\]When we substitute $n = 2022$, we get 
 \[f(2022) = \left(e^{\frac{2\pi i}{3}}\right)^{2022} + \left(e^{\frac{4\pi i}{3}}\right)^{2022}.\]We can rewrite $2022$ as $3 \cdot 674$, how does that help?
 \[f(2022) = \left(e^{\frac{2\pi i}{3}}\right)^{3 \cdot 674} + \left(e^{\frac{4\pi i}{3}}\right)^{3 \cdot 674} = \left(\left(e^{\frac{2\pi i}{3}}\right)^{3}\right)^{674} + \left(\left(e^{\frac{4\pi i}{3}}\right)^{3}\right)^{674} = 1^{674} + 1^{674} = \boxed{\textbf{(E)} \ 2}.\]Since any third root of unity must cube to $1$.

 



\textbf{Solution 2 (Eisenstein Units)}

The numbers $\frac{-1+i\sqrt{3}}{2}$ and $\frac{-1-i\sqrt{3}}{2}$ are both $\textbf{Eisenstein Units}$ (along with $1$), denoted as $\omega$ and $\omega^2$, respectively. They have the property that when they are cubed, they equal to $1$. Thus, we can immediately solve:
 \[\omega^{2022} + \omega^{2 \cdot 2022} = \omega^{3 \cdot 674} + \omega^{3 \cdot 2 \cdot 674} = 1^{674} + 1^{2 \cdot 674} = \boxed{\textbf{(E)} \ 2}.\]~mathboy100

 



\textbf{Solution 3 (Quick and Easy)}

We begin by recognizing this form looks similar to the definition of cosine: 
 \[\cos(x)=\frac{e^{ix}+e^{-ix}}{2}.\] We can convert our two terms into exponential form to find 
 \[f(n) = \left( e^{\frac{2\pi i}{3}} \right ) ^n + \left ( e^{-\frac{2\pi i}{3}} \right ) ^n=e^{\frac{2 \pi i n}{3}} + e^{-\frac{2\pi i n}{3}}.\] This simplifies nicely: 
 \[f(n)=2\cos\left( \frac{2\pi n}{3} \right).\] Thus, 
 \[f(2022)=2\cos \left ( \frac{2\pi (2022) }{3} \right) = 2\cos(1348 \pi) = \boxed{\textbf{(E)}\ 2}.\]




\textbf{Solution 4 (Third-order Homogeneous Linear Recurrence Relation)}

Notice how this looks like the closed form of the Fibonacci sequence except different roots. This is motivation to turn this closed formula into a recurrence relation. The base of the exponents are the roots of the characteristic equation $r^3-1=0$. So we have 
 \begin{align*} a_0&=2\\ a_1&=-1\\ a_2&=-1\\ a_n&=a_{n-3} \end{align*}
Every time $n$ is multiple of $3$ as is true when $n=2022$, $a_n= \boxed{\textbf{(E)} \ 2}$~lopkiloinm

 



\textbf{Solution 5 (Polynomial + Recursion)}

Let $a = \frac{-1+i\sqrt{3}}{2}$ and $b = \frac{-1-i\sqrt{3}}{2}$.We know that $a + b = -1$ and $a \cdot b = 1$.Therefore, a and b are the roots of $x^2 + x + 1 = 0$.By the factor theorem, $a^2 + a + 1 = 0$ and $b^2 + b + 1 = 0$.Multiply the first equation by $a^{n-2}$ and the second equation by $b^{n-2}$.This gives us $a^n + a^{n-1} + a^{n-2} = 0$ and $b^n + b^{n-1} + b^{n-2} = 0$.Adding both equations together we get $a^n + b^n + a^{n-1} + b^{n-1}+ a^{n-2} + b^{n-2} = 0$.This is the same as $f(n) + f(n-1) + f(n-2) = 0$.Therefore, $f(n) = -f(n-1) - f(n-2)$.Plugging in $n=1,2,3,4,5$, and $6$, we get $f(n) = -1, -1, 2, -1, -1, 2$. Therefore we know that if $n$ is a multiple of $3$, then $f(n)$ is $2$.Since $2022$ is a multiple of $3$, our answers is $E) 2$.~vpeddi18

 



\textbf{Solution 6 (SO FAST)}

Converting the two terms into rectangular form, 

 
 \[f(2022)=\left(\cos{\frac{2\pi}{3}}+i\sin{\frac{2\pi}{3}}\right)^{2022}+\left(\cos{\frac{4\pi}{3}}+i\sin{\frac{4\pi}{3}}\right)^{2022}.\]

 By DeMoivre's Theorem, 

 
 \[f(2022)=\left(\cos{\left(\frac{2\pi}{3}\cdot{2022}\right)}+i\sin{\left(\frac{2\pi}{3}\cdot{2022}\right)}\right)+\left(\cos{\left(\frac{4\pi}{3}\cdot{2022}\right)}+i\sin{\left(\frac{4\pi}{3}\cdot{2022}\right)}\right).\]

 Note that $\cos{\pi\cdot{k}}=1$ if $k$ is even and $-1$ if $k$ is odd, and that $\sin{\pi\cdot{k}}=0$ for all integers $k$. 

 All arguments are even in the second equation for $f(2022)$, so the two $\cos$ terms are equal to $1$, and the two $\sin$ terms are equal to $0$. 

 Therefore the answer is $1+1=\boxed{\textbf{(E) } 2}.$

 -Benedict T (countmath1)

 



\textbf{Solution 7 (No Euler's Identity)}

Observe that $(\frac{-1+i\sqrt{3}}{2})^{2}=\frac{-1-i\sqrt{3}}{2}$ and $(\frac{-1-i\sqrt{3}}{2})^{2}=\frac{-1+i\sqrt{3}}{2}$. 

 It follows that $(\frac{-1+i\sqrt{3}}{2})^{4}=\frac{-1+i\sqrt{3}}{2}$ and $(\frac{-1-i\sqrt{3}}{2})^{4}=\frac{-1-i\sqrt{3}}{2}$. 

 The binary expansion of 2022 is 
 \[2^{10}+2^9+2^8+2^7+2^6+2^5+2^2+2^1\] which can be rewritten as 
 \[(4^5+4^4+4^3+4^1)+2(4^4+4^3+4^2+4^0)\] Thus, 
 \[(\frac{-1+i\sqrt{3}}{2})^{2022}=((\frac{-1+i\sqrt{3}}{2})^4)^4\cdot (((\frac{-1+i\sqrt{3}}{2})^4)^4)^2\] or 
 \[(\frac{-1+i\sqrt{3}}{2})(\frac{-1-i\sqrt{3}}{2})\] Similarly, 
 \[(\frac{-1-i\sqrt{3}}{2})^{2022}=(\frac{-1-i\sqrt{3}}{2})(\frac{-1+i\sqrt{3}}{2})\] so 
 \[f(2022)=2(\frac{-1+i\sqrt{3}}{2})(\frac{-1-i\sqrt{3}}{2})=\boxed{\textbf{(E) } 2}\]



	\item (\textit{AMC 12 12B 2022}) Kayla rolls four fair $6$-sided dice. What is the probability that at least one of the numbers Kayla rolls is greater than $4$ and at least two of the numbers she rolls are greater than $2$?


 $\textbf{(A) }\frac{2}{3} \qquad \textbf{(B) }\frac{19}{27} \qquad \textbf{(C) }\frac{59}{81} \qquad \textbf{(D) }\frac{61}{81} \qquad \textbf{(E) }\frac{7}{9}$

 \textbf{Answer: }D



\textbf{Solution 1 (Complementary Counting)}

We will subtract from one the probability that the first condition is violated and the probability that \textit{only} the second condition is violated, being careful not to double-count the probability that both conditions are violated.

 For the first condition to be violated, all four dice must read $4$ or less, which happens with probability $\left( \frac23 \right)^4 = \frac{16}{81}$.

 For the first condition to be met but the second condition to be violated, at least one of the dice must read greater than $4$, but less than two of the dice can read greater than $2$. Therefore, one of the four die must read $5$ or $6$, while the remaining three dice must read $2$ or less, which happens with probability ${4 \choose 1} \left(\frac13\right) \left(\frac13\right)^3 = 4 \cdot \frac13 \cdot \frac{1}{27} = \frac{4}{81}$.

 Therefore, the overall probability of meeting both conditions is $1 - \frac{16}{81} - \frac{4}{81} = \boxed{\textbf{(D)}\ \frac{61}{81}}$.

 



\textbf{Solution 2 (Direct and Complementary Counting)}

There are either $0$, $1$, $2$, $3$, or $4$ dice that have values of more than $4$. The probability of getting $0$ is $\left(\frac{2}{3}\right)^4 = \frac{16}{81}$, the probability of getting $1$ is $4 \cdot \left(\frac{1}{3}\right)\left(\frac{2}{3}\right)^3 = \frac{32}{81}$, and the probability of getting $2$ or greater is $1 - \frac{16}{81} - \frac{32}{81} = \frac{11}{27}$. 

 It is obvious that the probability of getting at least two numbers greater than $2$ is $1$ if we have $2$ numbers greater than $4$.

 Let us calculate the probability of getting at least two numbers greater than $2$ if one die is greater than $4$ by using complementary counting. We already have one die that is greater than $2$, and the probability that a die that is less than $5$ is greater than $2$ is $\frac{1}{2}$. Thus, our probability is $1 - \left(\frac{1}{2}\right)^3 = \frac{7}{8}$.

 Finally, our total probability is $\frac{11}{27} + \frac{7}{8}\cdot\frac{32}{81} = \boxed{\textbf{(D)}\ \frac{61}{81}}$




\textbf{Solution 3 (Complementary Counting and PIE)}

We use Solution 1's idea of complementary counting. However, we will find the probability that the first condition is violated and the probability that the second condition is violated, then subtract the overlap (first and second condition both violated). We can do this because of PIE.

 As in Solution 1, the probability that the $\textbf{first condition}$ is violated is $\left(\frac{2}{3}\right)^4 = \frac{16}{81}.$

 The probability that the $\textbf{second condition}$ is violated (regardless of the first condition) can be broken into two cases.

 $\textbf{Case 1: exactly one number rolled is greater than 2}$

 Then, there are four choices for which die is the one greater than 2, $\frac{2}{3}$ chance for whether it reads 3, 4, 5, or 6, and $\left(\frac{1}{3}\right)^3 = \frac{1}{27}$ for the other three die that must read 1 or 2. The probability is then $4\left(\frac{2}{3}\right)\left(\frac{1}{27}\right) = \frac{8}{81}.$
$\textbf{Case 2: exactly zero numbers rolled are greater than 2}$

 Then, all four dice are either 1 or 2, so the probability of this is $\left(\frac{1}{3}\right)^4 = \frac{1}{81}.$

 Now, we look at the $\textbf{overlap}$. In other words, we want to find the probability none of the rolls are greater than $4$ and less than two of the rolls are greater than $2$. We split into the same two cases again.

 $\textbf{Case 1: exactly one number rolled is greater than 2 (and none of the rolls are greater than 4)}$

 There are still 4 choices for this die, and it must be 3 or 4, with probability $\frac{1}{3}$. The other must read 1 or 2, with probability $\frac{1}{27},$ so this case yields $4\left(\frac{1}{3}\right)\left(\frac{1}{27}\right) = \frac{4}{81}.$
$\textbf{Case 2: exactly zero numbers rolled are greater than 2 (and none of the rolls are greater than 4)}$

 The first statement implies the second statement, so this case yields $\frac{1}{81}$ as from our previous Case 2.

 Now we've got everything! Our answer is $1 - (\frac{16}{81} + (\frac{8}{81} + \frac{1}{81}) - (\frac{4}{81} + \frac{1}{81})) = \boxed{\textbf{(D)}\ \frac{61}{81}}.$




\textbf{Solution 4 (Direct Counting with Subtraction)}

Let X be the number of rolls greater than 4 (i.e. 5 or 6), and Y be the number of rolls greater than 2 (i.e. 3, 4, 5, or 6). We want to compute $P(X >= 1 and Y >= 2)$.

 We first compute $P(X >= 1)$. The complement is that no roll is greater than 4, so every roll is in {1,2,3,4}. The probability of this is $(4/6)^4 = 16/81$. Thus,

 $P(X >= 1) = 1 - 16/81 = 65/81$.

 This already guarantees Y >= 1, but we still need to enforce Y >= 2. Therefore, we subtract the outcomes where X >= 1 but Y = 1.

 The only way for X >= 1 and Y = 1 to occur is if exactly one die is greater than 2, and that die is also greater than 4. In other words, one die is 5 or 6, and the other three dice are 1 or 2.

 There are 4 choices for which die is >=5, and that die has probability 2/6 = 1/3. The other three dice must be in {1,2}, each with probability 1/3. Therefore, the probability of this “bad” case is

 $4 * (1/3) * (1/3)^3 = 4/81$.

 Finally, subtracting gives

 $P(X >= 1 and Y >= 2) = 65/81 - 4/81 = 61/81$.

 Thus, the answer is $\boxed{\textbf{(D)}\ 61/81}$.~FastOne

 


	\item (\textit{AMC 12 12B 2022}) The diagram below shows a rectangle with side lengths $4$ and $8$ and a square with side length $5$. Three vertices of the square lie on three different sides of the rectangle, as shown. What is the area of the region inside both the square and the rectangle?


 
                    \begin{figure}[H]
                        \centering
                        \includegraphics[width=0.5\textwidth]{Problem/AMC_12_12B_2022_13_Problem_Figure_1.png}
                    \end{figure}
                    
$\textbf{(A) }15\dfrac{1}{8}  \qquad \textbf{(B) }15\dfrac{3}{8}  \qquad \textbf{(C) }15\dfrac{1}{2}  \qquad \textbf{(D) }15\dfrac{5}{8}  \qquad \textbf{(E) }15\dfrac{7}{8}$

 \textbf{Answer: }D



\textbf{Solution 1}

Let us label the points on the diagram.

 
                    \begin{figure}[H]
                        \centering
                        \includegraphics[width=0.5\textwidth]{Solution/AMC_12_12B_2022_13_Solution_Figure_1.png}
                    \end{figure}
                    

 By doing some angle chasing using the fact that $\angle ACE$ and $\angle CEG$ are right angles, we find that $\angle BAC = \angle DCE = \angle FEG$. Similarly, $\angle ACB = \angle CED = \angle EGF$. Therefore, $\triangle ABC \sim \triangle CDE \sim \triangle EFG$.

 As we are given a rectangle and a square, $AB = 4$ and $AC = 5$. Therefore, $\triangle ABC$ is a $3$-$4$-$5$ right triangle and $BC = 3$. 

 $CE$ is also $5$. So, using the similar triangles, $CD = 4$ and $DE = 3$.

 $EF = DF - DE = 4 - 3 = 1$. Using the similar triangles again, $EF$ is $\frac14$ of the corresponding $AB$. So, 

 
 \begin{align*} [\triangle EFG] &= \left(\frac14\right)^2 \cdot [\triangle ABC] \\ &= \frac{1}{16} \cdot 6 \\ &= \frac38. \end{align*}


 Finally, we have 

 
 \begin{align*} [ACEG] &= [ABDF] - [\triangle ABC] - [\triangle CDE] - [\triangle EFG] \\ &= 7 \cdot 4 - \frac12 \cdot 3 \cdot 4 - \frac12 \cdot 3 \cdot 4 - \frac38 \\ &= 28 - 6 - 6 - \frac38 \\ &= \boxed{\textbf{(D) }15\dfrac{5}{8}}. \end{align*}





\textbf{Solution 2 (Clever)}

(Refer to the diagram above)Proceed the same way as Solution 1 until you get all of the side lengths. Then, it is clear that due to the answer choices, we only need to find the fractional part of the shaded area. The area of the whole rectangle is integral, as is the area of $\triangle ABC$, $\triangle CDE$, and the rectangle to the far left of the diagram. The area of $EFG$ is $\frac{3}{8}$ and thus the fractional part of the answer is $\frac{5}{8}$. The only answer choice that has $\frac{5}{8}$ in it is $\boxed{\textbf{(D) }15\dfrac{5}{8}}$




\textbf{Solution 3 (Intuitive Risking)}


                    \begin{figure}[H]
                        \centering
                        \includegraphics[width=0.5\textwidth]{Solution/AMC_12_12B_2022_13_Solution_Figure_2.png}
                    \end{figure}
                    We can quickly find out through Pythagorean Triples that \( GH = 3 \). Now, we make an educated guess that \( G \) is the midpoint of \( FI \), and therefore \( HI = 1 \).

 Noice! So, now what? Well, \( FG \) is then \( 4 \), and \( EF \) would then be \( 3 \).

 Given that info, we also see that \( DE = 1 \).

 Our area of the rectangle \( ADFI \) is \( 8 \times 4 = 32 \). We then see that we must subtract rectangle \( ABHI = 4 \times 1 = 4 \), triangle \( BHG = \frac{1}{2} \times 3 \times 4 = 6 \), and triangle \( GFE = \frac{1}{2} \times 4 \times 3 = 6 \). This gives \( 32 - (4 + 6 + 6) = 16 \).

 We make a guess that the intersection between \( B \) and \( D \), call it \( J \), to \( E \) is exactly \( 5/4 \). Then, through the Pythagorean theorem, \( JD = 1/2 \), and the area is \( \frac{1}{2} \times \frac{1}{2} \times 1 = \frac{1}{4} \). So, the area of that mini triangle is \( \frac{1}{4} \), in reality we know it is bigger. We subtract this from \( 16 \) to get \( 15 \frac{3}{4} \) and take the closest answer choice that is lower than the result we got, which is $\boxed{\textbf{(D) }15\dfrac{5}{8}}$.

 Take risks everyone! You might do great things like solve an AMC problem :D!

 



\textbf{Solution 4 (Coordinate Geometry)}

Same diagram as Solution 1, but added point $H$, which is $(4,7)$. I also renamed all the points to form coordinates using $B$ as the origin. 
                    \begin{figure}[H]
                        \centering
                        \includegraphics[width=0.5\textwidth]{Solution/AMC_12_12B_2022_13_Solution_Figure_3.png}
                    \end{figure}
                    

 In order to find the area, point $G$'s coordinates must be found. Notice how $EH$ and $AG$ intercept at point $G$. This means that we need to find the equations for $EH$ and $AG$ and make a system of linear equations. 

 Using the slope formula $m=\frac{y_{2} - y_{1}}{x_{2} - x_{1}}$, we get the slope for $EH$, which means $m=\frac{3-7}{7-4} = -\frac{4}{3}$

 Then, by using point slope form. $y-y_{1}=m(x-x_{1})$. We can say that the equation for $EH$ is $y-7=-\frac{4}{3}(x-4)$ or in this case, $y=-\frac{4}{3}x+12 \frac{1}{3}$.

 And it is easy to figure out that the equation for $AG$ is $y=4$.

 The best way to solve the system of linear equations is to substitute the $y$ for the $4$ in equation $EH$. $4=-\frac{4}{3}x+12 \frac{1}{3}$, so $x=6\frac{1}{4}$ and $y=4$ This would mean $G\left(6\frac{1}{4},4\right)$.

 Since we have our $G$ coordinate, we can continue with Solution 3, with the area of the trapezoid $\left(\frac{EG+AC}{2}\right)(CE)$, where $EG=\frac{5}{4}$ (using distance formula for $E$ to $G$), $AC=5$, and $CE=5$.

 By substitution, we get $\left(\frac{\frac{5}{4}+5}{2}\right)(5)=$$\boxed{\textbf{(D) }15\dfrac{5}{8}}$. 

 



\textbf{Solution 5 (FAST)}

Notice the small triangle in the upper right corner is a $3-4-5$ triangle. Then that triangle is similar to the big triangle by AA similarity. From that, do similar triangles and you find that the longer leg of the small triangle is 1. Then you find that the triangle below is $3-4-5$, so the side length of the rectangle (without the outer rectangle) is 7. afterwards you just add the half of the square + the remaining triangle which can be found by multiplying base and height (in which we already know)

 
I found the same solution when solving this, but essentially, we know all triangles in the picture are similar, and we know the lower left hand triangle is a $3-4-5$ triangle, and it is similar to the top triangle, where 4 corresponds to 5, therefore the unknown side length of that triangle is 15/4. Continuing on, the are of the white part of the square is $5 * 15/4 * 1/2 = 75/8,  25 - 75/8 = \fbox{(D)}$

 
-kaiser

 



\textbf{Solution 6 (Fastest Similar Triangles)}

For reference, use the points labelled in the diagram of Solution 3. Let the point one unit to the right of $A$ be $A'$ (so that $A'$ is one of the vertices of the square). The square means $A'C = 5$, so we get a $3$-$4$-$5$ triangle $A'BC$.

 $m\angle BA'C = 90 ^{\circ} - m\angle CA'G = m\angle GA'H$. 

 Therefore $\triangle GA'H$ is proportional to a $3$-$4$-$5$ triangle, with $HG$ corresponding to $3$ and $A'H$ corresponding to $4$. By similar triangles, we find 

 $HG = A'H \cdot \frac{3}{4} = \frac{15}{4}$.

 
            \begin{enumerate}[label=(\roman*)]
                <dt>then, finally, $[A'CEG] = [A'CEH] - [\triangle GA'H] = 5^2 - \frac{1}{2} \cdot \frac{15}{4} \cdot 5 = \boxed{\textbf{(D) }15\dfrac{5}{8}}.$</dt>
            \end{enumerate}
            


	\item (\textit{AMC 12 12B 2022}) The graph of $y=x^2+2x-15$ intersects the $x$-axis at points $A$ and $C$ and the $y$-axis at point $B$. What is $\tan(\angle ABC)$?


 $\textbf{(A) }\frac{1}{7} \qquad \textbf{(B) }\frac{1}{4} \qquad \textbf{(C) }\frac{3}{7} \qquad \textbf{(D) }\frac{1}{2} \qquad \textbf{(E) }\frac{4}{7}$

 \textbf{Answer: }E



\textbf{Solution 1 (Dot Product)}

First, find $A=(-5,0)$, $B=(0,-15)$, and $C=(3,0)$. Create vectors $\overrightarrow{BA}$ and $\overrightarrow{BC}.$ These can be reduced to $\langle -1, 3 \rangle$ and $\langle 1, 5 \rangle$, respectively. Then, we can use the dot product to calculate the cosine of the angle (where $\theta=\angle ABC$) between them:

 
 \begin{align*} \langle -1, 3 \rangle \cdot \langle 1, 5 \rangle = 15-1 &= \sqrt{10}\sqrt{26}\cos(\theta),\\ \implies \cos (\theta) &= \frac{7}{\sqrt{65}}. \end{align*}


 Thus, 
 \[\tan(\angle ABC) = \sqrt{\frac{65}{49}-1}= \boxed{\textbf{(E)}\ \frac{4}{7}}.\]




\textbf{Solution 2}

Note that $y=x^2+2x-15$ intersects the $x$-axis at points $(-5, 0)$ and $(3, 0)$. Without loss of generality, let these points be $A$ and $C$ respectively. Also, the graph intersects the $y$-axis at point $B = (0, -15)$.

 Let point $O=(0, 0)$. It follows that $\triangle AOB$ and $\triangle BOC$ are right triangles. 

 We have 
 \[\tan(\angle ABC) = \tan(\angle ABO + \angle OBC) = \frac{\tan(\angle ABO) + \tan(\angle OBC)}{1 - \tan(\angle ABO) \cdot \tan(\angle OBC)} = \frac{\frac15 + \frac13}{1 - \frac1{15}} = \boxed{\textbf{(E)}\ \frac{4}{7}}.\] Alternatively, we can use the \href{https://artofproblemsolving.com/wiki/index.php?title=Pythagorean\_Theorem}{Pythagorean Theorem} to find that $AB = 5 \sqrt{10}$ and $BC = 3 \sqrt{26}$ and then use the $A = \frac12 ab \sin \angle C$ \href{https://artofproblemsolving.com/wiki/index.php?title=Area\_of\_a\_Triangle}{area formula for a triangle} and the \href{https://artofproblemsolving.com/wiki/index.php?title=Law\_of\_Cosines}{Law of Cosines} to find $\tan(\angle ABC)$.

 



\textbf{Solution 3}

Like above, we set $A$ to $(-5,0)$, $B$ to $(0, -15)$, and $C$ to $(3,0)$, then finding via the Pythagorean Theorem that $AB = 5 \sqrt{10}$ and $BC = 3 \sqrt{26}$. Using the Law of Cosines, we see that 
 \[\cos(\angle ABC) = \frac{AB^2 + BC^2 - AC^2}{2 AB BC} = \frac{250 + 234 - 64}{30 \sqrt{260}} = \frac{7}{\sqrt{65}}.\] Then, we use the identity $\tan^2(x) = \sec^2(x) - 1$ to get 
 \[\tan(\angle ABC) = \sqrt{\frac{65}{49} - 1} = \boxed{\textbf{(E)}\ \frac{4}{7}}.\]

 Alternatively, we could observe that $\sin(\angle ABC)=\sqrt{1-\cos^2(\angle ABC)}=\sqrt{1-\left(\dfrac7{\sqrt{65}}\right)^2}=\sqrt{\dfrac{16}{65}}=\dfrac4{\sqrt{65}}$, so $\tan(\angle ABC)=\dfrac{\sin(\angle ABC)}{\cos(\angle ABC)}=\boxed{\textbf{(E) }\dfrac47}$. 

 



\textbf{Solution 4}

We can reflect the figure, but still have the same angle. This problem is the same as having points $D(0,0)$, $E(3,15)$, and $F(-5,15)$, where we're solving for angle FED. We can use the formula for $\tan(a-b)$ to solve now where $a$ is the $x$-axis to angle $F$ and $b$ is the $x$-axis to angle $E$. $\tan(a) = \textrm{slope of line }DF = -3$ and $\tan(B) = \textrm{slope of line }DE = 5$. Plugging these values into the $\tan(a-b)$ formula, we get $(-3-5)/(1+(-3\cdot 5))$ which is $\boxed{\textbf{(E)}\ \frac{4}{7}}.$




\textbf{Solution 5}

We use the formula $[ABC]=\frac{1}{2}ab\sin{C}.$

 Note that $\triangle ABC$ has side-lengths $AB=5\sqrt{10}$ and $BC=3\sqrt{26}$ from Pythagorean theorem, with the area being $\frac12\cdot8\cdot15.$

 We equate the areas together to get: 
 \[\frac12\cdot8\cdot15=\frac12\cdot5\sqrt{10}\cdot3\sqrt{26}\cdot\sin{B},\]from which $\sin{B}=\frac{8}{\sqrt{260}}.$

 From Pythagorean Identity, $\cos{B}=\frac{14}{\sqrt{260}}.$

 Then we use $\tan{B}=\frac{\sin{B}}{\cos{B}}$, to obtain $\tan{B}=\frac{8}{14}=\boxed{\textbf{(E)}\ \frac{4}{7}}.$

 - SAHANWIJETUNGA

 



\textbf{Solution 6 (Complex Numbers)}


                    \begin{figure}[H]
                        \centering
                        \includegraphics[width=0.5\textwidth]{Solution/AMC_12_12B_2022_14_Solution_Figure_1.png}
                    \end{figure}
                    From $x^2 + 2x - 15 = (x-3)(x+5)$, we may assume, without loss of generality, that $x$-intercepts of the given parabola are $A(-5,0)$ and $C(3,0)$. And, point $B$ has coordinates $(0,-15)$. Consider complex numbers $z = 3 + i$ and $w = 5 + i$ whose arguments are $\theta \coloneqq \angle OBA$ and $\phi \coloneqq \angle OBC$, respectively. Notice that $\angle ABC = \theta + \phi$ is the argument of the product $zw$ which is 
 \[zw = (3+i)(5+i) = 14 + 8i.\]Hence 
 \[\tan \angle ABC = \frac{\operatorname{Im}(zw)}{\operatorname{Re}(zw)} = \frac{8}{14} = \boxed{\textbf{(E)}\ \frac{4}{7}}.\]



	\item (\textit{AMC 12 12B 2022}) One of the following numbers is not divisible by any prime number less than $10.$ Which is it?


 $\textbf{(A) } 2^{606}-1 \qquad\textbf{(B) } 2^{606}+1 \qquad\textbf{(C) } 2^{607}-1 \qquad\textbf{(D) } 2^{607}+1\qquad\textbf{(E) } 2^{607}+3^{607}$

 \textbf{Answer: }C



\textbf{Solution 1 (Modular Arithmetic)}

For $\textbf{(A)}$ modulo $3,$
 \begin{align*} 2^{606} - 1 & \equiv (-1)^{606} - 1 \\ & \equiv 1 - 1 \\ & \equiv 0 . \end{align*}
Thus, $2^{606} - 1$ is divisible by $3.$

 For $\textbf{(B)}$ modulo $5,$
 \begin{align*} 2^{606} + 1 & \equiv 2^{{\textrm Rem} ( 606, \phi(5) )} + 1 \\ & \equiv 2^{{\textrm Rem} ( 606, 4 )} + 1 \\ & \equiv 2^2 + 1 \\ & \equiv 0 . \end{align*}
Thus, $2^{606} + 1$ is divisible by $5.$

 For $\textbf{(D)}$ modulo $3,$
 \begin{align*} 2^{607} + 1 & \equiv (-1)^{607} + 1 \\ & \equiv - 1 + 1 \\ & \equiv 0 . \end{align*}
Thus, $2^{607} + 1$ is divisible by $3.$

 For $\textbf{(E)}$ modulo $5,$
 \begin{align*} 2^{607} + 3^{607} & \equiv 2^{607} + (-2)^{607} \\ & \equiv 2^{607} - 2^{607} \\ & \equiv 0 . \end{align*}
Thus, $2^{607} + 3^{607}$ is divisible by $5.$

 Therefore, the answer is $\boxed{\textbf{(C) }2^{607} - 1}.$




\textbf{Solution 2 (Factoring)}

We have
 \begin{alignat*}{8} 2\textasciicircum{}{606}-1 &= 4\textasciicircum{}{303}-1 &&= (4-1)(4\textasciicircum{}{302}+4\textasciicircum{}{301}+4\textasciicircum{}{300}+\cdots+4\textasciicircum{}0), \\  2\textasciicircum{}{606}+1 &= 4\textasciicircum{}{303}+1 &&= (4+1)(4\textasciicircum{}{302}-4\textasciicircum{}{301}+4\textasciicircum{}{300}-\cdots+4\textasciicircum{}0), \\  2\textasciicircum{}{607}+1 & &&= (2+1)(2\textasciicircum{}{606}-2\textasciicircum{}{605}+2\textasciicircum{}{604}-\cdots+2\textasciicircum{}0), \\ 2\textasciicircum{}{607}+3\textasciicircum{}{607} & &&= (2+3)(2\textasciicircum{}{606}\cdot3\textasciicircum{}0-2\textasciicircum{}{605}\cdot3\textasciicircum{}1+2\textasciicircum{}{604}\cdot3\textasciicircum{}2-\cdots+2\textasciicircum{}0\cdot3\textasciicircum{}{606}). \end{alignat*}We conclude that $\textbf{(A)}$ is divisible by $3$, $\textbf{(B)}$ is divisible by $5$,$\textbf{(D)}$ is divisible by $3$, and $\textbf{(E)}$ is divisible by $5$.

 Since all of the other choices have been eliminated, we are left with $\boxed{\textbf{(C) } 2^{607}-1}$.

 



\textbf{Solution 3 (Elimination)}

Mersenne Primes are primes of the form $2^n-1$, where $n$ is prime. Using the process of elimination, we can eliminate every option except for $\textbf{(A)}$ and $\textbf{(C)}$. Clearly, $606$ isn't prime, so the answer must be $\boxed{\textbf{(C) }2^{607}-1}$.

 Note: A organization of computer scientist and mathematicians named Great Internet Mersenne Prime Search (GIMPS) search for the worlds biggest prime numbers. More information can be found on \href{https://www.mersenne.org/primes/}{https://www.mersenne.org/primes/}~hashbrown2009

 



\textbf{Solution 3a (Elimination)}

We examine option E first. $2^{607}$ has a units digit of $8$ (Taking the units digit of the first few powers of two gives a pattern of $2, 4, 8, 6, 2, 4, 8, 6, 2, 4, 8, 6,\cdots$) and $3^{607}$ has a units digit of $7$ (Taking the units digit of the first few powers of three gives a pattern of $3, 9, 7, 1, 3, 9, 7, 1, 3, 9, 7, 1,\cdots$). Adding $7$ and $8$ together, we get $15$, which is a multiple of $5$, meaning that $2^{607}+3^{607}$ is divisible by 5.

 Next, we examine option D. We take the first few powers of $2$ added with $1$:
 \[2^1+1=3\]
 \[2^2+1=5\]
 \[2^3+1=9\]
 \[2^4+1=17\]
 \[2^5+1=33\]
 \[2^6+1=65\]
 \[2^7+1=129\]

 We see that the odd powers of $2$ added with 1 are multiples of three. If we continue this pattern, $2^{607}+1$ will be divisible by $3$. (The reason why this pattern works: When you multiply $2 \equiv2\pmod{3}$ by $2$, you obtain $4 \equiv1 \pmod{3}$. Multiplying by $2$ again, we get $1\cdot2\equiv2 \pmod{3}$. We see that in every cycle of two powers of $2$, it goes from $2 \pmod{3}$ to $1 \pmod{3}$ and back to $2 \pmod{3}$.)

 Next, we examine option B. We see that $2^{606}$ has a units of digits of $4$ (Taking the units digit of the first few powers of two gives a pattern of $2, 4, 8, 6, 2, 4, 8, 6, 2, 4, 8, 6,\cdots$). Adding $1$ to $4$, we get $5$. Since $2^{606}+1$ has a units digit of $5$, it is divisible by $5$.

 Lastly, we examine option A. Using the difference of cubes factorization $a^3-b^3=(a-b)(a^2+ab+b^2)$, we have $2^{606}-1^3=(2^{202}-1)(2^{404}+2^{202}+1)$. Since $2^{404}+2^{202}+1\equiv0\pmod{3}$ (Every term in the sequence is equivalent to $1\pmod{3}$), $2^{606}-1$ is divisible by $3$.

 Since we have eliminated every option except C, $\boxed{\text{(C)} \hspace{0.1 in}2^{607}-1}$ is not divisible by any prime less than $10$.

 



\textbf{Solution 3b (Elimination + Number Theory)}

We know that the prime numbers less than 10 are $2,3,5$ and $7$. We can start by testing if any of the answer choices are divisible by $2$. We see that they are all sums of one even number and one odd number, which is simply odd. So, we cannot exclude any answer choices so far. Now, let's check divisibility by $3$. We can use the fact that $2 \equiv -1 \pmod{3}$ to our advantage: 

 


 $(\text{A})\hspace{0.1in} 2^{606}-1$$\equiv (-1)^{606} -1$$\equiv 1-1$ $\equiv 0 \pmod{3}$
$(\text{B}) \hspace{0.1in} 2^{606}+1$$\equiv (-1)^{606} +1$$\equiv 1+1$$\equiv 2 \pmod{3}$
$(\text{C})\hspace{0.1in} 2^{607}-1$$\equiv (-1)^{607} -1$$\equiv -1 - 1$$\equiv -2 \pmod{3}$
$(\text{D})\hspace{0.1in} 2^{607} +1 \equiv (-1)^{607} +1 \equiv -1+1 \equiv 0 \pmod{3}$
$(\text{E}) \hspace{0.1in}  2^{607} + 3^{607} \equiv 2^{607} \equiv (-1)^{607} \equiv -1 \pmod{3}$

 
So, we eliminate choices A and D from divisibility by 3. Now, we move onto divisibility by 5. We can use cycling of powers to find useful remainders. Let's start with choice $B$. 

 


 $(\text{B}) \hspace{0.1in} 2^{606}+1$ 

 


 We see that the remainders of powers of $2$ when divided by $5$ cycle in a pattern: $2,4,3,1,2,4,3,1 \dots$. Since the pattern cycles every 4 terms, we use modulo 4 to simplify 606, getting that $606 \equiv 2\pmod{4}$. So, we get that $2^{606}$ is congruent modulo 5 to the second term of our pattern, so $2^{606} \equiv{4} \pmod{5}$. Thus, $2^{606}+1 \equiv 4+1 \equiv 5\equiv 0\pmod{5}$. We now eliminate choice B and compare choices C and E.  

 
$(\text{E}) \hspace{0.1in} 2^{607} + 3^{607}$ 

 


 
Looking at choice E, we see that we have to do similar cycling for powers of $3$. We get a pattern of $3,4,2,1,3,4,2,1 \dots$. Since this pattern also cycles every 4 terms, we use modulo 4 to simplify 607, getting that $607 \equiv 3\pmod{4}$. Both the power of 3 and the power of 2 have an exponent of 607, so we use the third term (since we just found that $607 \equiv 3\pmod{4}$) in each corresponding 4 term pattern to get that $2^{607} + 3^{607} \equiv 3+2 \equiv 5 \equiv 0\pmod{5}$. We eliminate choice E, and we are left with the correct answer: choice $\boxed{(\text{C}) \hspace{0.1in} 2^{607}-1}$




\textbf{Solution 3c (Elimination, Number Theory, and Function Analysis)}

Let’s start by calculating the unit’s digit of each answer choice.

 $2^{606}$ has a unit’s digit of 4, so choices A and B have unit’s digits of $3$ and $5$ respectively.

 $2^{607}$ has a unit’s digit of 8, so choices C and D have unit’s digits of $7$ and $9$ respectively.

 $3^{607}$ has a unit’s digit of 7. Combining with the statement above, choice E has a unit’s digit of $5$.

 So, none of the answer choices are divisible by $2$. But, we can eliminate choices B and E because they both have unit’s digit of $5$, and therefore, are divisible by $5$.

 We are left now with choices A, C, and D. Let’s now apply function analysis.

 Look at choices A and C again. They share the form $2^{n}-1$. Notice that if $n$ is even, it follows that $2^{n}-1$ is divisible by $3$. Since $606$ is even, choice A is divisible by $3$. However, our answer cannot be divisible by $3$, so eliminate A. This leaves us with choices C and D.

 Look at choice D again. It has the form $2^{n}+1$. Notice that if $n$ is odd, it follows that $2^{n}+1$ is divisible by 3. Since $607$ is odd, choice D is divisible by $3$. However, since the answer cannot be divisible by $3$, we eliminate D.

 At this point, you would choose choice C, but let’s now test divisibility by $7$. Notice that if $n$ is a multiple of 3, it follows that $2^{n}-1$ is divisible by $7$. However, since $607$ is not a multiple of $3$, choice C is not divisible by $7$.

 Since it is not divisible by any prime less than $10$, $\boxed{(\text{C}) \hspace{0.1in} 2^{607}-1}$ is the correct answer.

 



\textbf{Solution 4 (modulo 15 elimination)}

We proceed by taking the solutions modulo $15$, allowing us to eliminate both multiples of $3$ and multiples of $5$.

 First, we must identify cycles of powers of 2 modulo 15.$2^0 \equiv 1, \quad 2^1 \equiv 2, \quad 2^2 \equiv 4, \quad 2^3 \equiv 8, \quad 2^4 \equiv 1 \pmod{15}.$Thus, the pattern repeats every $4$ powers:$2^n \equiv \begin{cases} 1 &\text{if } n \equiv 0 \pmod{4},\\ 2 &\text{if } n \equiv 1 \pmod{4},\\ 4 &\text{if } n \equiv 2 \pmod{4},\\ 8 &\text{if } n \equiv 3 \pmod{4}. \end{cases}$

 Since $606 \equiv 2 \pmod{4}$ and $607 \equiv 3 \pmod{4}$,$2^{606} \equiv 4 \pmod{15}, \qquad 2^{607} \equiv 8 \pmod{15}.$

 Now let's look at cycles of powers of 3 mod 15.$3^0 \equiv 1, \quad 3^1 \equiv 3, \quad 3^2 \equiv 9, \quad 3^3 \equiv 12, \quad 3^4 \equiv 6, \quad 3^5 \equiv 3 \pmod{15}.$The cycle from here repeats every $4$ terms for $n \ge 1$:$3^n \equiv \begin{cases} 6 &\text{if } n \equiv 0 \pmod{4},\\ 3 &\text{if } n \equiv 1 \pmod{4},\\ 9 &\text{if } n \equiv 2 \pmod{4},\\ 12 &\text{if } n \equiv 3 \pmod{4}. \end{cases}$

 Since $607 \equiv 3 \pmod{4}$,$3^{607} \equiv 12 \pmod{15}$.

 We then evaluate each option modulo 15.

 \begin{aligned}\textbf{(A)} &:\; 2\textasciicircum{}{606}-1 \equiv 4-1 \equiv 3 \pmod{15},\\[6pt]\textbf{(B)} &:\; 2\textasciicircum{}{606}+1 \equiv 4+1 \equiv 5 \pmod{15},\\[6pt]\textbf{(C)} &:\; 2\textasciicircum{}{607}-1 \equiv 8-1 \equiv 7 \pmod{15},\\[6pt]\textbf{(D)} &:\; 2\textasciicircum{}{607}+1 \equiv 8+1 \equiv 9 \pmod{15},\\[6pt]\textbf{(E)} &:\; 2\textasciicircum{}{607}+3\textasciicircum{}{607} \equiv 8+12 \equiv 20 \equiv 5 \pmod{15}.\end{aligned}

 The residues divisible by $3$ or $5$ modulo $15$ are 0, 3, 5, 6, 9, 10, 12.Among our results, only $7$ (from option $\boxed{(\text{C})}$) is not among these, eliminating all the other answer choices.

 
$\boxed{(\text{C}) \hspace{0.1in} 2^{607}-1}$ is the number not divisible by any prime less than $10$. \\~Loquacious\_Autodidact

 



\textbf{Solution 5 (Contradiction; Less than 1 minute if you're good at number theory)}

This is like solution 1, but don't find the remainder, we assume it.

 The primes less than 10 are 2, 3, 5, and 7.

 Option choice $A$: Let $2^{606} \equiv 1 \pmod{3}$ as it clearly isn't divisible by 2. Then, because $2^2 \equiv 1 \pmod{3}$, it follows that $2^{606}$ satisfies the modulo. Option $A$ is not possible.

 Option choice $B$ ends in 5, and thus not possible.

 Option choice $C$: It is obviously not divisible by 2. We set $2^{607} \equiv 1 \pmod{3}$, which we see is not possible through Fermat's Little Theorem. We see $2^{607}$ doesn't end in 5. We set $2^{607} \equiv 1 \pmod{7}$, and through Fermat's Little Theorem $2^6 \equiv 1 \pmod{7}$. It follows that $607$ isn't a factor of 2, and thus does not hold the desired claim, a contradiction.

 At this point, for speed purposes, one would circle $C$ and move on. However, given the directory that this is a solution, I shall show the last two cases are not possible.

 Option choice $D$: It is obviously not divisible by 2. We set $2^{607} \equiv 2 \pmod{3}$. Through properties of modulo 3 we have that $(-1)^{607} \equiv -1 \pmod{3}$, which is true. Thus, $D$ is not plausible.

 Option choice $E$ is divisible by 5, and thus not possible. $\square$

 Additionally, for option $D$. One could notice that $2^{607} \equiv 2 \pmod{3}$ in the following way. We analyze the pattern. $2^1 \equiv 2 \pmod{3}$. $2^3 \equiv 2 \pmod{3}$. $2^5 \equiv 2 \pmod{3}$. $2^{2n+1} \pmod{3} \equiv 4^n \cdot 2^1 \pmod{3} \equiv 1^n \cdot 2 \pmod{3} \equiv 2 \pmod{3}$. 607 is odd, and thus $2^{607} \equiv 2 \pmod{3}$, as desired. $\square$




\textbf{Solution 6 (Extremely Illegal Cheese)}

All of the other solutions are extremely long, but this is how to cheese it using some methods that MIGHT be unreliable:First, turn 606 to 0 and 607 to 1. Next, Calculate each value, giving you $A=0 B=2 C=1 D=3 E=5$. Now since 2, 3, and 5 are prime, they can be divided by a prime(themselves). 0 doesn't really have multiples, so you can literally ignore it, and since 1 is only divisible by 1, which isnt prime, you get the correct answer, $\boxed{(\text{C})}$.

 


	\item (\textit{AMC 12 12B 2022}) Suppose $x$ and $y$ are positive real numbers such that

 \[x^y=2^{64}\text{ and }(\log_2{x})^{\log_2{y}}=2^{7}.\]
What is the greatest possible value of $\log_2{y}$?


 $\textbf{(A) }3 \qquad \textbf{(B) }4 \qquad \textbf{(C) }3+\sqrt{2} \qquad \textbf{(D) }4+\sqrt{3} \qquad \textbf{(E) }7$

 \textbf{Answer: }C



\textbf{Solution 1}

Take the base-two logarithm of both equations to get
 \[y\log_2 x = 64\quad\text{and}\quad (\log_2 y)(\log_2\log_2 x) = 7.\]Now taking the base-two logarithm of the first equation again yields
 \[\log_2 y + \log_2\log_2 x = 6.\]It follows that the real numbers $r:=\log_2 y$ and $s:=\log_2\log_2 x$ satisfy $r+s=6$ and $rs = 7$. Solving this system yields
 \[\{\log_2 y,\log_2\log_2 x\}\in\{3-\sqrt 2, 3 + \sqrt 2\}.\]Thus the largest possible value of $\log_2 y$ is $3+\sqrt 2 \implies \boxed{\textbf{(C) }3+\sqrt{2}}$.

 



\textbf{Solution 2}

$x^y=2^{64} \Rightarrow y\log_2{x}=64 \Rightarrow \log_2{x}=\dfrac{64}{y}$.

 Substitution into $(\log_2{x})^{\log_2{y}}=2^{7}$ yields

 $\left(\dfrac{64}{y}\right)^{\log_2{y}}=2^{7} \Rightarrow \log_2{y}\log_2{\dfrac{64}{y}}=7 \Rightarrow \log_2{y}(6-\log_2{y})=7 \Rightarrow \log^2_2{y}-6\log_2{y}+7=0$.

 Solving for $\log_2{y}$ yields $\log_2{y}=3-\sqrt{2}$ or $3+\sqrt{2}$, and we take the greater value $\boxed{\textbf{(C) }3+\sqrt{2}}$.

 



\textbf{Solution 3}

Let $x = 2^a, y = 2^b.$ We have $(2^a)^{2^b} = 2^{64} \Rightarrow 2^{a\cdot 2^b} = 2^{64} \Rightarrow a\cdot 2^b = 64,$ and $a^b = 128$.

 Then, from eq 1, $a = 64\cdot 2^{-b},$ and substituting in to eq 2, $(64\cdot 2^{-b})^b = 64^b\cdot 2^{-b^2} = 2^{6b}\cdot 2^{-b^2} = 2^{6b-b^2} = 2^{7}.$ Thus, $6b-b^2 = 7.$

 Solving for $b$ using the quadratic formula gets $b = 3 \pm \sqrt{2}.$ Since we are looking for $\log_2{y}$ which equals $b,$ we put $\boxed{\textbf{(C) }3+\sqrt{2}}$ as our answer.

 


	\item (\textit{AMC 12 12B 2022}) How many $4 \times 4$ arrays whose entries are $0$s and $1$s are there such that the row sums (the sum of the entries in each row) are $1, 2, 3,$ and $4,$ in some order, and the column sums (the sum of the entries in each column) are also $1, 2, 3,$ and $4,$ in some order? For example, the array

 \[\left[   \begin{array}{cccc}     1 & 1 & 1 & 0 \\     0 & 1 & 1 & 0 \\     1 & 1 & 1 & 1 \\     0 & 1 & 0 & 0 \\   \end{array} \right]\]
satisfies the condition.


 $\textbf{(A) }144 \qquad \textbf{(B) }240 \qquad \textbf{(C) }336 \qquad \textbf{(D) }576 \qquad \textbf{(E) }624$

 \textbf{Answer: }D



\textbf{Solution 1 (One-to-One Correspondence)}

Note that the arrays and the sum configurations have one-to-one correspondence. Furthermore, the row sum configuration and the column sum configuration are independent of each other. Therefore, the answer is $(4!)^2=\boxed{\textbf{(D) }576}.$
<u>\textbf{Remark}</u>

 For any given sum configuration, we can uniquely reconstruct the array it represents. Conversely, for any array, it is clear that we can determine the unique sum configuration associated with it. Therefore, this establishes a one-to-one correspondence between the arrays and the sum configurations.

 



\textbf{Solution 2 (Linear Transformation and Permutation)}

In this problem, we call a matrix that satisfies all constraints given in the problem a feasible matrix.

 First, we observe that if a matrix is feasible, and we swap two rows or two columns to get a new matrix, then this new matrix is still feasible.

 Therefore, any feasible matrix can be obtained through a sequence of such swapping operations from a feasible matrix where for all 

 


	\item (\textit{AMC 12 12B 2022}) Each square in a $5 \times 5$ grid is either filled or empty, and has up to eight adjacent neighboring squares, where neighboring squares share either a side or a corner. The grid is transformed by the following rules:


 
            \begin{itemize}
        \item Any filled square with two or three filled neighbors remains filled.
            \end{itemize}
            
            \begin{itemize}
        \item Any empty square with exactly three filled neighbors becomes a filled square.
            \end{itemize}
            
            \begin{itemize}
        \item All other squares remain empty or become empty.
            \end{itemize}
            A sample transformation is shown in the figure below.

                    \begin{figure}[H]
                        \centering
                        \includegraphics[width=0.5\textwidth]{Problem/AMC_12_12B_2022_18_Problem_Figure_1.png}
                    \end{figure}
                    
Suppose the $5 \times 5$ grid has a border of empty squares surrounding a $3 \times 3$ subgrid. How many initial configurations will lead to a transformed grid consisting of a single filled square in the center after a single transformation? (Rotations and reflections of the same configuration are considered different.)

                    \begin{figure}[H]
                        \centering
                        \includegraphics[width=0.5\textwidth]{Problem/AMC_12_12B_2022_18_Problem_Figure_2.png}
                    \end{figure}
                    
$\textbf{(A)}\ 14 \qquad\textbf{(B)}\ 18 \qquad\textbf{(C)}\ 22 \qquad\textbf{(D)}\ 26 \qquad\textbf{(E)}\ 30$

 \textbf{Answer: }C



\textbf{Solution}

There are two cases for the initial configuration:

 
            \begin{enumerate} 
        \item The center square is filled. Exactly two of the eight adjacent neighboring squares of the center are filled. Clearly, the only possibility is that the squares along one diagonal are filled, as shown below:

 
                    \begin{figure}[H]
                        \centering
                        \includegraphics[width=0.5\textwidth]{Solution/AMC_12_12B_2022_18_Solution_Figure_1.png}
                    \end{figure}
                    In this case, there are $2$ possible initial configurations. All rotations and reflections are considered.

  

 
        \item The center square is empty.Exactly three of the eight adjacent neighboring squares of the center are filled. The possibilities are shown below:

 
                    \begin{figure}[H]
                        \centering
                        \includegraphics[width=0.5\textwidth]{Solution/AMC_12_12B_2022_18_Solution_Figure_2.png}
                    \end{figure}
                    In this case, there are $4+4+4+8=20$ possible initial configurations. All rotations and reflections are considered.

 

 
            \end{enumerate}
            Together, the answer is $2+20=\boxed{\textbf{(C)}\ 22}.$




\textbf{Solution Alternate}

Key Idea:

 Only the $3$ x $3$ inner grid can initially contain filled squares, since the outer border is empty. We want exactly the center square to be filled after one step.

 Case 1: The center is initially filled.

 A filled square remains filled precisely when it has $2$ or $3$ filled neighbors. We examine the possible choices of filled neighbors of the center and require that no other square becomes filled.

 This gives $2$ valid configurations, counting rotations/reflections.

 Case 2: The center is initially empty.

 For the center to become filled, it must have exactly $3$ filled neighbors. There are several possible shapes for these three neighbors. Counting the configurations of each shape gives

 $2 + 4 + 4 + 4 + 8 = 22.$

 The total number of configurations is

 \begin{align*}2+(2+4+4+4+8)&=2+22\\&=\boxed{24}.\end{align*}


 Therefore the answer is $\boxed{\textbf{(C)}\ 22}.$

 \href{https://artofproblemsolving.com/wiki/index.php?title=User:Lunar\_paradox&action=edit&redlink=1}{Lunar paradox} (\href{https://artofproblemsolving.com/wiki/index.php?title=User\_talk:Lunar\_paradox&action=edit&redlink=1}{talk}) 18:06, 24 August 2026 (EDT)

 \textbf{Remark} This problem is based off of Conway's Game of Life, a cellular automaton created by British mathematician John Conway in 1970. The game follows the exact same pattern of rules as this problem, but is played on an infinite grid of cells.

 


	\item (\textit{AMC 12 12B 2022}) In $\triangle{ABC}$ medians $\overline{AD}$ and $\overline{BE}$ intersect at $G$ and $\triangle{AGE}$ is equilateral. Then $\cos(C)$ can be written as $\frac{m\sqrt p}n$, where $m$ and $n$ are relatively prime positive integers and $p$ is a positive integer not divisible by the square of any prime. What is $m+n+p?$
$\textbf{(A) }44 \qquad \textbf{(B) }48 \qquad \textbf{(C) }52 \qquad \textbf{(D) }56 \qquad \textbf{(E) }60$

 \textbf{Answer: }A



\textbf{Solution 1 (Law of Cosines)}

Let $AG=AE=EG=2x$. Since $E$ is the midpoint of $\overline{AC}$, we must have $EC=2x$.

 Since the centroid splits the median in a $2:1$ ratio, $GD=x$ and $BG=4x$.

 Applying Law of Cosines on $\triangle ADC$ and $\triangle{}AGB$ yields $AB=\sqrt{28}x$ and $CD=BD=\sqrt{13}x$. Finally, applying Law of Cosines on $\triangle ABC$ yields $\cos(C)=\frac{5}{2\sqrt{13}}=\frac{5\sqrt{13}}{26}$. The requested sum is $5+13+26=44$.

 



\textbf{Solution 2 (Law of Cosines: One Fewer Step)}

Let $AG = 1$. Since $\frac{BG}{GE}=2$ (as $G$ is the centroid), $BE = 3$. Also, $EC = 1$ and $\angle{BEC} = 120^{\circ}$. By the law of cosines (applied on $\triangle BEC$), $BC = \sqrt{13}$.

 Applying the law of cosines again on $\triangle BEC$ gives $\cos{\angle{C}} = \frac{1 + 13 - 9}{2\sqrt{13}} = \frac{5\sqrt{13}}{26}$, so the answer is $\fbox{\textbf{(A)}\ 44}$.

 



\textbf{Solution 3 (Law of Cosine)}

Let $AG = AE = GE = CE = 1$. Since $G$ is the centroid, $DG = \frac12$, $BG = 2$.

 
 \[\angle BGD = \angle AGE = 60^{\circ}\]

 By the Law of Cosine in $\triangle BGD$

 \[BD^2 = BG^2 + DG^2 - 2 \cdot BG \cdot DG \cdot \cos \angle BGD\]

 \[BD = \sqrt {2^2 + \left( \frac{1}{2} \right)^2 - 2 \cdot 2 \cdot \frac12 \cdot \cos \angle BGD} = \frac{\sqrt{13}}{2}, \quad CD = \frac{\sqrt{13}}{2}\]

 By the Law of Cosine in $\triangle ACD$

 \[AD^2 = AC^2 + CD^2 - 2 \cdot AC \cdot CD \cdot \cos \angle C\]

 \[\cos \angle C = \frac{ AC^2 + CD^2 - AD^2 }{ 2 \cdot AC \cdot CD } = \frac{ 2^2 + \left( \frac{\sqrt{13}}{2} \right)^2 - \left( \frac{3}{2} \right)^2 }{ 2 \cdot 2 \cdot \frac{\sqrt{13}}{2} } = \frac{ 5 \sqrt{13} }{26}\]

 \[5 + 13 + 26 = \boxed{\textbf{(A) }44}\]




\textbf{Solution 4 (Barycentric Coordinates)}

Using reference triangle $\triangle AGE$, we can let 
 \[A=(1,0,0),G=(0,1,0),E=(0,0,1),C=(-1,0,2),D=(-\frac{1}{2},\frac{3}{2},0),B=(0,3,-2).\] If we move $A,B,C$ each over by $(1,0,-2)$, leaving $\angle C$ unchanged, we have 
 \[A=(2,0,-2),B=(1,3,-4),C=(0,0,0).\] The angle $\theta$ between vectors $\overrightarrow{CA}$ and $\overrightarrow{CB}$ satisfies 
 \[\cos\theta=\frac{(2)(1)+(0)(3)+(-2)(-4)}{\sqrt{\left[2^{2}+0^{2}+(-2)^{2}\right]\left[1^{2}+3^{2}+(-4)^{2}\right]}}=\frac{10}{\sqrt{8\cdot 26}}=\frac{10}{4\sqrt{13}}=\frac{5\sqrt{13}}{26},\] giving the answer, $5+13+26=\boxed{\textbf{(A)}~44}$.

 



\textbf{Solution 5 (Coordinate Bash)}

Let \(A\) be at \((0,0)\), with \(E\) at \((1,0)\) and \(G\) at \((\frac{1}{2},\frac{\sqrt{3}}{2})\). Because \(\triangle AGE\) is equilateral, the equation of line \(AD\) is \(y=x\sqrt{3}\), and the equation for \(BE\) is \(y=-x\sqrt{3}+\sqrt{3}\). Because \(BE\) is a median, we know that \(C\) is at \((2,0)\).  

 Therefore, the equation for line \(BC\) is \(y=m(x-2)\), where \(m\) is the slope of the line. Since \(AD\) is also a median, \(D\) is the midpoint of \(\overline{BC}\).  

 To find \(m\), solve for the coordinates of \(D\) in terms of \(m\):  

 
 \[x\sqrt{3}=mx-2m\] 
 \[x(m-\sqrt{3})=2m\] 
 \[x=\frac{2m}{m-\sqrt{3}}\] 

 The \(y\)-value of \(D\) is then \(y=x\sqrt{3}\):  

 
 \[y=\frac{2m\sqrt{3}}{m-\sqrt{3}}\] 

 Because \(D\) is the midpoint of \(\overline{BC}\) and \(y_C=0\), the \(y\)-value of \(B\) is double that of \(D\):  

 
 \[y_B=\frac{4m\sqrt{3}}{m-\sqrt{3}}\] 

 Now use the equation for line \(BE\) to express \(B\):  

 
 \[mx-2m=-x\sqrt{3}+\sqrt{3}\] 
 \[x(m+\sqrt{3})=2m+\sqrt{3}\] 
 \[x_B=\frac{2m+\sqrt{3}}{m+\sqrt{3}}\] 

 Plug this back into \(BE\) to find \(y_B\):  

 
 \[y_B=\frac{-2m-\sqrt{3}}{m+\sqrt{3}}+\sqrt{3}=\frac{-m\sqrt{3}}{m+\sqrt{3}}\] 

 Equating the two expressions for \(y_B\):  

 
 \[-\frac{m\sqrt{3}}{m+\sqrt{3}}=\frac{4m\sqrt{3}}{m-\sqrt{3}}\] 
 \[-\sqrt{3}m^2+3m=4\sqrt{3}m^2+12m\] 
 \[5\sqrt{3}m^2+9m=0\] 
 \[m(5m\sqrt{3}+9)=0\] 

 Since \(m\neq0\), we have  

 
 \[m=-\frac{3\sqrt{3}}{5}.\] 

 Now find \(B\):  

 
 \[-x\sqrt{3}+\sqrt{3}=-\frac{3\sqrt{3}}{5}x+\frac{6\sqrt{3}}{5}\] 
 \[\frac{-2\sqrt{3}}{5}x=\frac{\sqrt{3}}{5}\] 
 \[x_B=-\frac{1}{2}\] 

 Plug this into \(BE\) to get  

 
 \[y_B=\frac{3\sqrt{3}}{2}.\] 

 Since we’re asked to find \(\cos\angle C\), extend \(\overline{AC}\) to \(F\) where \(F=(-\frac{1}{2},0)\), forming right triangle \(\triangle BFC\). Using the Pythagorean theorem for \(\overline{BC}\):  

 
 \[\left(\frac{5}{2}\right)^2+\left(\frac{3\sqrt{3}}{2}\right)^2=BC^2,\quad BC=\sqrt{13}.\] 

 Thus  

 
 \[\cos\angle C=\frac{\frac{5}{2}}{\sqrt{13}}=\frac{5\sqrt{13}}{26}.\] 

 The problem asks for \(m+n+p\), so \(5+13+26=\boxed{44}\). 

 


	\item (\textit{AMC 12 12B 2022}) Let $P(x)$ be a polynomial with rational coefficients such that when $P(x)$ is divided by the polynomial $x^2 + x + 1$, the remainder is $x + 2$, and when $P(x)$ is divided by the polynomial $x^2 + 1$, the remainder is $2x + 1$. There is a unique polynomial of least degree with these two properties. What is the sum of the squares of the coefficients of that polynomial?


 $\textbf{(A) } 10 \qquad \textbf{(B) } 13 \qquad \textbf{(C) } 19 \qquad \textbf{(D) } 20 \qquad \textbf{(E) } 23$

 \textbf{Answer: }E



\textbf{Solution 1 (Experimentation)}

Given that all the answer choices and coefficients are integers, we hope that $P(x)$ has positive integer coefficients.

 Throughout this solution, we will express all polynomials in base $x$. E.g. $x^2 + x + 1 = 111_{x}$.

 We are given:
 \[111a + 12 = 101b + 21 = P(x).\]We add $111$ and $101$ to each side and balance respectively:
 \[111(a - 1) + 123 = 101(b - 1) + 122 = P(x).\]We make the unit's digits equal:
 \[111(a - 1) + 123 = 101(b - 2) + 223 = P(x).\]We now notice that:
 \[111(a - 11) + 1233 = 101(b - 12) + 1233 = P(x).\]Therefore $a = 11_{x} = x + 1$, $b = 12_{x} = x + 2$, and $P(x) = 1233_{x} = x^3 + 2x^2 + 3x + 3$. $3$ is the minimal degree of $P(x)$ since there is no way to influence the $x$‘s digit in $101b + 21$ when $b$ is an integer. The desired sum is $1^2 + 2^2 +3^2+ 3^2 = \boxed{\textbf{(E)} \ 23}$

 P.S. The four computational steps can be deduced through quick experimentation.

 



\textbf{Solution 2}

Let $P(x) = Q(x)(x^2+x+1) + x + 2$, then $P(x) = Q(x)(x^2+1) + xQ(x) + x + 2$, therefore $xQ(x) + x + 2 \equiv 2x + 1 \pmod{x^2+1}$, or $xQ(x) \equiv x-1 \pmod{x^2+1}$. The rule says we can change any $x^2$ into $-1$. This is because $x^2 \equiv -1 \pmod{x^2+1}$ (since $x^2 - (-1) = x^2+1$, which is a multiple of $x^2+1$).To find the least possible degree $Q(x)$, it would be logical to first assume it's a constant. Let's say $Q(x)=c$, which makes $xQ(x) = cx$. Is $cx$ congruent to $x-1$? The difference is $cx - (x-1) = (c-1)x+1$. This is a polynomial of degree 1, so it's not a multiple of $x^2+1$ (a degree 2 polynomial). So, a constant won't work.Could $Q(x)$ be a simple degree 1 polynomial? Let's try $Q(x) = x+1$. Let's check if $Q(x) = x+1$ works. Calculate $xQ(x)$: $x(x+1) = x^2+x$. 

 Now, apply our rule: we can replace $x^2$ with $-1$. Therefore, $x^2+x$ becomes $-1+x$, which is $x-1$. This shows that $xQ(x)$ is congruent to $x-1$.Since we showed that no constant polynomial works, and we found a degree 1 polynomial that does work, the minimum degree must be 1. The simplest polynomial of degree 1 that satisfies the condition is $Q(x)=x+1$, and expanding gives $P(x) = x^3+2x^2+3x+3$. Summing the squares of coefficients gives $\boxed{\textbf{(E)} \ 23}$




\textbf{Solution 3}

Let $P(x) = (x^2+x+1)Q_1(x) + x + 2$, then $P(x) = (x^2+1)Q_1(x) + xQ_1(x) + x + 2$

 Also, $P(x) = (x^2+1)Q_2(x) + 2x + 1$

 We infer that $Q_1(x)$ and $Q_2(x)$ have same degree, we can assume $Q_1(x) = x + a$, and $Q_2(x) = x + b$, since $P(x)$ has least degree. If this cannot work, we will try quadratic, etc.

 Then we get:$(x^2+1)(Q_1(x) - Q_2(x)) + xQ_1(x) - x + 1 = 0$

 The constant term gives us:$(Q_1(x) - Q_2(x)) + 1 = 0$

 So $Q_1(x) - Q_2(x) = -1$

 Substituting this in gives: $-(x^2+1) + xQ_1(x) - x + 1 = 0$

 Solving this equation, we get$Q_1(x) = x + 1$ 

 Plugging this into our original equation we get $P(x) = x^3 + 2x^2 + 3x + 3$

 Verify this works with $P(x) = (x^2+1)Q_2(x) + 2x + 1$

 Therefore the answer is $1^2 + 2^2 + 3^2 + 3^2 = \boxed{\textbf{(E)} \ 23}$




\textbf{Solution 4 (Undetermined Coefficients)}

Notice that we cannot have the quotients equal to some constants, since the same constant will yield different constant terms for $P(x)$ (which is bad) and different constants will yield different first coefficients (also bad). Thus, we try setting the quotients equal to linear terms (for minimizing degree).

 Let $P(x)=(x^2+x+1)(ax+b)+(x+2)$ and $P(x)=(x^2+1)(ax+c)+(2x+1)$. The quotients have the same $x$ coefficient, since $P(x)$ must have the same $x^3$ coefficient in both cases. Expanding, we get 
 \[P(x)=ax^3+(a+b)x^2+(a+b+1)x+(b+2)\] and 
 \[P(x)=ax^3+cx^2+(a+2)x+(c+1).\]

 Equating coefficients, we get $b+2=c+1$, $a+b+1=a+2$, and $a+b=c$. From the second equation, we get $b=1$, then substituting into         the first, $c=2$. Finally, from $a+b=c$, we have $a=1$. Now, $P(x)=(x^2+x+1)(ax+b)+(x+2)=(x^2+x+1)(x+1)+(x+2)=x^3+2x^2+3x+3$ and our answer is 
 \[1^2+2^2+3^2+3^2=\boxed{\textbf{(E)} \ 23}.\]




\textbf{Solution 5 (Quick, but Not Quicker Than Solution 2)}

We construct the following equations in terms of $P(x)$ and the information given by the problem:
 \[\textbf{(1) } P(x)=(x^2+x+1)\cdot Q(x)+x+2\]
 \[\textbf{(2) } P(x)=(x^2+1)\cdot R(x)+2x+1\]Upon inspection, $Q(x)$ and $R(x)$ cannot be constant, so the smallest possible degree of $P(x)$ is $3,$ and both $Q(x)$ and $R(x)$ are linear. 

 Let $Q(x)=x-q$ and $R(x)=x-r.$ We know there will be values for $q$ and $r$ that make the below equation hold, so we can assume that $P(x)$ has a leading coefficient of $1$.

 Substituting these values in, and setting $\textbf{(1)}$ and $\textbf{(2)}$ equal to each  other,
 \[(x^2+x+1)(x-q)+x+2=(x^2+1)(x-r)+2x+1.\]We plug in $x=0$, yielding $r+1=q.$ Substituting this value into the above equation, 
 \[(x^2+x+1)(x-r-1)+x+2=(x^2+1)(x-r)+2x+1.\]Letting $x=1,$ we conclude that $r=-2,$ so $R(x)=x+2.$Therefore, 
 \[P(x)=(x^2+1)(x+2)+2x+1 = x^3+2x^3+3x+3.\]The requested sum is 
 \[1^2+2^2+3^2+3^2=\boxed{\textbf{(E) }23}\]

 -Benedict T (countmath1)

 



\textbf{Solution 6 (Similar to Solution 3)}

By remainder theorem, the polynomial can be written as follows.

 
 \[P(x) = (x^2+x+1)Q_{1}(x)+x+2 = (x^2+1)Q_{2}(x)+2x+1\]This is a timed exam, we can use the information given by answer choices. The answer choices tell us this is the polynomial with integer coefficients, and we need to find the polynomial with the least degree so we can assume both $Q_{1}(x)$ and $Q_{2}(x)$ are linear (the coefficient of x should be same).

 Then we can write $P(x)$ as a cubic polynomial.

 
 \[P(x) = (x^2+x+1)(ax+b)+x+2 = (x^2+1)(ax+c)+2x+1\]Substituting $x=0,1,-1$ to determine the value of $a$ and $b$.

 We have:
 \[P(0) = b+2 = c+1\]
 \[P(1) = 3a+3b+3 = 2a+2c+3\]
 \[P(-1) = -a+b+1 = -2a+2c-1\]

 We can solve the simultaneous equations: $a=1,b=1,c=2$

 Hence, $P(x)=(x^2+x+1)(x+1)+x+2=x^3+2x^2+3x+3$. The answer is $1^2+2^2+3^2+3^2=\boxed{\textbf{(E) }23}$




\textbf{Solution 7 (Bounding with Answer Choices)}

We are given: 

 $P(x) \equiv x + 2 \pmod{x^2+x+1}$
$P(x) \equiv 2x + 1 \pmod{x^2+1}$

 Since the problem asks us about the coefficients of the polynomial, we should evaluate $P(x)$ at $1$.

 Our congruences now become:

 $P(1) \equiv 1 + 2 \pmod{(1)^2+(1)+1}$ $\implies P(1) \equiv 3 \pmod{3} \implies P(1) \equiv 0 \pmod{3}$.

 $P(1) \equiv 2(1) + 1 \pmod{(1)^2+1}$ $\implies P(1) \equiv 3 \pmod{2} \implies P(1) \equiv 1 \pmod{2}$.

 We now know that the coefficients of the polynomial sum to an odd multiple of 3. 

 We now begin bounding the value of the sum of the coefficients squared. For our answer, the sum of the coefficients squared must be between $10$ and $23$. Trying a quadratic polynomial, if $P(1) = 3$, then the maximum of the coefficients squared is $9$ ($3^2+0^2+0^2$ is the maximum, any other permutation can be trivially shown to be lower). If $P(1) = 9$, the minimum of the coefficients squared is $27$ ($3^2+3^2+3^2$, once again, any other permutation can be trivially shown to be higher). Any higher value of $P(1)$ obviously will not work, so $P(x)$ cannot be a quadratic.

 Since a quadratic did not work, we now move up to a cubic. We start with $P(1) = 3$. Similarly to the quadratic, there is no way to get an answer within the range of the answer choices. 

 However, moving up to $P(1) = 9$, we get that the minimum of the coefficients squared is $81/4$ (which we get from all 4 coefficients being equal at $9/4$). Since this is the minimum, any other answers will be higher. We get that $\boxed{\textbf{(E) }23}$ is our answer, since it is the only answer choice greater than $81/4$.

 -SwordOfJustice

 



\textbf{Solution 8}

Since the divisors $x^2 + x + 1$ and $x^2 + 1$ are both powers of 2, and the remainders $x + 2$ and $2x + 1$ are both linear, we can assume that $P(x)$ has a minimum degree of 3.

 We can then set the coefficients of the polynomial as $P(x) = ax^3 + bx^2 + cx + d$.

 Then, using long division, we divide  $P(x) = ax^3 + bx^2 + cx + d$ by $x^2 + x + 1$ and obtain a remainder of $(c - b)x + (a + d - b)$, which is equal to $x + 2$

 We also divide $P(x) = ax^3 + bx^2 + cx + d$ by $x^2 + 1$ and obtain a remainder of $(c - a)x + (d - b)$, which is equal to $2x + 1$.

 From here we can set up a systems of equations:

 $c -b = 1$
$a + d - b = 2$
$c - a = 2$
$d - b = 1$

 Thus giving us the solutions $a = 1$, $b = 2$, $c = 3$, $d = 3$ and hence $P(x) = x^3 + 2x^2 + 3x + 3$

 Squaring the coefficients, then adding them, $1^2 + 2^2 + 3^2 + 3^2$ we get the final answer $\boxed{\textbf{(E)} \ 23}$




\textbf{Solution 9}

We assume $P(x)$ is a monic cubic polynomial. Now, we set $P(x)=x^{3}+ax^{2}+bx+c$, and set the equations:

 $x^{3}+ax^{2}+bx+c=Q_{1}(x)(x^{2}+x+1)+x+2$
$x^{3}+ax^{2}+bx+c=Q_{2}(x)(x^{2}+1)+2x+1$

 Plugging in $x=i$ for the second equation, we get $-i-a+bi+c=2i+1$, and matching coefficients gives $c-a=1$ and $b-1=2 \rightarrow b=3$. Now, we can guess values for $a$ and $c$ based on the answer choices. We have that the sum of the squares of the current known coefficients is $1^{2}+3^{2}=10$, and guessing $c=1, a=0$ clearly doesnt match up with the answer choices, so we guess $c=3, a=2$ and find that the answer is $\boxed{\textbf{(E)} \ 23}$




\textbf{Solution 10 (EZ CRT)}

If we let $x=10$, we just get a CRT problem. Note: this only works because all the coefficients are integers between 0 and 9. Our problem now is to find $n$ such that
 \[n\equiv 12\pmod {111}\]
 \[n\equiv 21\pmod {101}\]After using some careful Euclidean Algorithm (or just eyeing it), $n=1233$, so our answer is $1^2+2^2+3^2+3^2 = \boxed{\textbf{(E)} \ 23}$.

 



\textbf{Solution 11 (Polynomial CRT)}

$x^2 + x + 1$ and $x^2 + 1$ are relatively prime according to Bézout's theorem.

 $P(x) \equiv U(x) + V(x) \pmod{(x^2 + x + 1)(x^2 + 1)}$, where
 \[U(x) \equiv x + 2 \pmod{x^2 + x + 1}, \quad U(x) \equiv 0 \pmod{x^2 + 1}\]
 \[V(x) \equiv 2x + 1 \pmod{x^2 + 1}, \quad V(x) \equiv 0 \pmod{x^2 + x + 1}\]

 Notice that\begin{align}& x^2 + 1 \equiv -x \pmod{x^2 + x + 1}\\& (x^2 + 1)(x + 1) \equiv -x(x + 1) \equiv 1 \pmod{x^2 + x + 1}\\& U(x) \equiv -(x^2 + 1) + 2(x^2 + 1)(x + 1) \pmod{(x^2 + x + 1)(x^2 + 1)}\\& U(x) \equiv 2x^3 + x^2 + 2x + 1 \pmod{(x^2 + x + 1)(x^2 + 1)}\end{align}

 And\begin{align}& x^2 + x + 1 \equiv x \pmod{x^2 + 1}\\& (x^2 + x + 1)x \equiv x^2 \equiv -1 \pmod{x^2 + 1}\\& V(x) \equiv 2(x^2 + x + 1) - (x^2 + x + 1)x \pmod{(x^2 + x + 1)(x^2 + 1)}\\& V(x) \equiv -x^3 + x^2 + x + 2 \pmod{(x^2 + x + 1)(x^2 + 1)}\end{align}

 Therefore $P(x) \equiv U(x) + V(x) \equiv x^3 + 2x^2 + 3x + 3 \pmod{(x^2 + x + 1)(x^2 + 1)}$, the answer is $1^2 + 2^2 + 3^2 + 3^2 = \boxed{\textbf{(E)} \ 23}$



	\item (\textit{AMC 12 12B 2022}) Let $S$ be the set of circles in the coordinate plane that are tangent to each of the three circles with equations $x^{2}+y^{2}=4$, $x^{2}+y^{2}=64$, and $(x-5)^{2}+y^{2}=3$. What is the sum of the areas of all circles in $S$?


 $\textbf{(A) } 48 \pi \qquad \textbf{(B) } 68 \pi \qquad \textbf{(C) } 96 \pi \qquad \textbf{(D) } 102 \pi \qquad \textbf{(E) } 136 \pi \qquad$

 \textbf{Answer: }E



\textbf{Solution 1}


                    \begin{figure}[H]
                        \centering
                        \includegraphics[width=0.5\textwidth]{Solution/AMC_12_12B_2022_21_Solution_Figure_1.png}
                    \end{figure}
                    The circles match up as follows: Case $1$ is brown, Case $2$ is blue, Case $3$ is green, and Case 4 is red.Let $x^2 + y^2 = 64$ be circle $O$, $x^2 + y^2 = 4$ be circle $P$, and $(x-5)^2 + y^2 = 3$ be circle $Q$.All the circles in S are internally tangent to circle $O$.There are four cases with two circles belonging to each:

 $*$ $P$ and $Q$ are internally tangent to $S$.

 $*$ $P$ and $Q$ are externally tangent to $S$.

 $*$ $P$ is externally and Circle $Q$ is internally tangent to $S$.

 $*$ $P$ is internally and Circle $Q$ is externally tangent to $S$.

 Consider Cases $1$ and $4$ together. Since circles $O$ and $P$ have the same center, the line connecting the center of $S$ and the center of $O$ will pass through the tangency point of both $S$ and $O$ and the tangency point of $S$ and $P$. This line will be the diameter of $S$ and have length $r_P + r_O = 10$. Therefore the radius of $S$ in these cases is $5$.

 Consider Cases $2$ and $3$ together. Similarly to Cases $1$ and $4$, the line connecting the center of $S$ to the center of $O$ will pass through the tangency points. This time, however, the diameter of $S$ will have length $r_P-r_O=6$. Therefore, the radius of $S$ in these cases is $3$.

 The set of circles $S$ consists of $8$ circles - $4$ of which have radius $5$ and $4$ of which have radius $3$. The total area of all circles in $S$ is $4(5^2\pi + 3^2\pi) = 136\pi \Rightarrow \boxed{\textbf{(E)}}$.

 -naman12

 



\textbf{Solution 2}

We denote by $C_1$ the circle that has the equation $x^2 + y^2 = 4$.We denote by $C_2$ the circle that has the equation $x^2 + y^2 = 64$.We denote by $C_3$ the circle that has the equation $(x-5)^2 + y^2 = 3$.

 We denote by $C_0$ a circle that is tangent to $C_1$, $C_2$ and $C_3$.We denote by $\left( u, v \right)$ the coordinates of circle $C_0$, and $r$ the radius of this circle.

 From the graphs of circles $C_1$, $C_2$, $C_3$, we observe that if $C_0$ is tangent to all of them, then $C_0$ must be internally tangent to $C_2$.We have
 \[ u^2 + v^2 = \left( 8 - r \right)^2 . \hspace{1cm} (1) \]

 We do the following casework analysis in terms of the whether $C_0$ is externally tangent to $C_1$ and $C_3$.

 Case 1: $C_0$ is externally tangent to $C_1$ and $C_3$.

 We have
 \[ u^2 + v^2 = \left( r + 2 \right)^2   \hspace{1cm} (2) \]and
 \[ (u-5)^2 + v^2 = \left( r + \sqrt{3} \right)^2 . \hspace{1cm} (3) \]

 Taking $(2) - (1)$, we get $r + 2 = 8 - r$. Thus, $r = 3$.We can further compute (omitted here) that there exist feasible $(u,v)$ with this given $r$.

 Case 2: $C_1$ is internally tangent to $C_0$ and $C_3$ is externally tangent to $C_0$.

 We have
 \[ u^2 + v^2 = \left( r - 2 \right)^2  \hspace{1cm} (2) \]and
 \[ (u-5)^2 + v^2 = \left( r + \sqrt{3} \right)^2 . \hspace{1cm} (3) \]

 Taking $(2) - (1)$, we get $r - 2 = 8 - r$. Thus, $r = 5$.We can further compute (omitted here) that there exist feasible $(u,v)$ with this given $r$.

 Case 3: $C_1$ is externally tangent to $C_0$ and $C_3$ is internally tangent to $C_0$.

 We have
 \[ u^2 + v^2 = \left( r + 2 \right)^2  \hspace{1cm} (2) \]and
 \[ (u-5)^2 + v^2 = \left( r - \sqrt{3} \right)^2 . \hspace{1cm} (3) \]

 Taking $(2) - (1)$, we get $r + 2 = 8 - r$. Thus, $r = 3$.We can further compute (omitted here) that there exist feasible $(u,v)$ with this given $r$.

 Case 4: $C_1$ is internally tangent to $C_0$ and $C_3$ is internally tangent to $C_0$.

 We have
 \[ u^2 + v^2 = \left( r - 2 \right)^2  \hspace{1cm} (2) \]and
 \[ (u-5)^2 + v^2 = \left( r - \sqrt{3} \right)^2 . \hspace{1cm} (3) \]

 Taking $(2) - (1)$, we get $r - 2 = 8 - r$. Thus, $r = 5$.We can further compute (omitted here) that there exist feasible $(u,v)$ with this given $r$.

 
Because the graph is symmetric with the $x$-axis, and for each case above, the solution of $v$ is not 0. Hence, in each case, there are two congruent circles whose centers are symmetric through the $x$-axis.

 Therefore, the sum of the areas of all the circles in $S$ is $2\left( 3^2 \pi +5^2 \pi +3^2 \pi +5^2 \pi \right) = \boxed{\textbf{(E) } 136 \pi}$.

 



\textbf{Solution 3}

The problem of counting the number of circles tangent to a given number of other circles, lines, or points is called Apollonius's Problem. Here, there are three circles; solutions to Apollonius's problem tell us that there are 8 possible tagent circles. When looking for all eight, we can notice that the diagram is symmetrical about the x-axis, so we only need to find four different circles and then reflect them to the other side. 

 After drawing in those circles, we can use any of the above methods to find the radii of the circle. The important step is to realize that the segment between tangent circles' centers must pass through the point of tangency. From here, we can prove that the center of the circle with equation $x^2 + y^2 = 4$ and the center of the circle tangent to all three must both lie on the radius of $x^2 + y^2 = 64$, allowing us to quickly find the radius of the tangent circle.

 For the tangent circles mostly above the x-axis, there are two of radius 5, and two of radius 3. The total area is $50 \pi + 18 \pi = 68 \pi$. As due to symmetry there are four other circles below the x-axis, the total area is $136 \pi$, so we choose $\boxed{ (E) }$. 

 


 


	\item (\textit{AMC 12 12B 2022}) Ant Amelia starts on the number line at $0$ and crawls in the following manner. For $n=1,2,3,$ Amelia chooses a time duration $t_n$ and an increment $x_n$ independently and uniformly at random from the interval $(0,1).$ During the $n$th step of the process, Amelia moves $x_n$ units in the positive direction, using up $t_n$ minutes. If the total elapsed time has exceeded $1$ minute during the $n$th step, she stops at the end of that step; otherwise, she continues with the next step, taking at most $3$ steps in all. What is the probability that Amelia’s position when she stops will be greater than $1$?


 $\textbf{(A) }\frac{1}{3} \qquad \textbf{(B) }\frac{1}{2} \qquad \textbf{(C) }\frac{2}{3} \qquad \textbf{(D) }\frac{3}{4} \qquad \textbf{(E) }\frac{5}{6}$

 \textbf{Answer: }C



\textbf{Solution 1 (Geometric Probability)}

Let $x$ and $y$ be random variables that are independently and uniformly distributed in the interval $(0,1).$ Note that 
 \[P(x+y\leq 1)=\frac{\frac12\cdot1^2}{1^2}=\frac12,\] as shown below:
                    \begin{figure}[H]
                        \centering
                        \includegraphics[width=0.5\textwidth]{Solution/AMC_12_12B_2022_22_Solution_Figure_1.png}
                    \end{figure}
                    Let $x,y,$ and $z$ be random variables that are independently and uniformly distributed in the interval $(0,1).$ Note that 
 \[P(x+y+z\leq 1)=\frac{\frac13\cdot\left(\frac12\cdot1^2\right)\cdot1}{1^3}=\frac16,\] as shown below:
                    \begin{figure}[H]
                        \centering
                        \includegraphics[width=0.5\textwidth]{Solution/AMC_12_12B_2022_22_Solution_Figure_2.png}
                    \end{figure}
                    We have two cases:

 
            \begin{enumerate} 
        \item Amelia takes exactly $2$ steps.We need $x_1+x_2>1$ and $t_1+t_2>1.$ So, the probability is 
 \[P(x_1+x_2>1)\cdot P(t_1+t_2>1)=\left(1-\frac12\right)\cdot\left(1-\frac12\right)=\frac14.\]

        \item Amelia takes exactly $3$ steps.We need $x_1+x_2+x_3>1$ and $t_1+t_2\leq1.$ So, the probability is 
 \[P(x_1+x_2+x_3>1)\cdot P(t_1+t_2\leq1)=\left(1-\frac16\right)\cdot\frac12=\frac{5}{12}.\]

            \end{enumerate}
            Together, the answer is $\frac14 + \frac{5}{12} = \boxed{\textbf{(C) }\frac{2}{3}}.$




\textbf{Solution 2 (Generalization and Induction)}

We can in fact find the probability that any number of randomly distributed numbers on the interval $[0, 1]$ sum to more than $1$ using geometric probability, as shown in the video below.

 If we graph the points that satisfy $x + y < 1$, $0 < x, y < 1$, we get the triangle with points $(0, 0)$, $(1, 0)$, and $(0, 1)$. If we graph the points that satisfy $x + y + z < 1$, $0 < x, y, z < 1$, we get the tetrahedron with points $(0, 0, 0)$, $(1, 0, 0)$, $(0, 1, 0)$, and $(0, 0, 1)$.

 Of course, the probability of either of these cases happening is simply the area/volume of the points we graphed divided by the total area of the graph, which is always $1$ (this would be much simpler than my calculus proof above).

 Thus, we can now solve for the probability that the sum is less than one for $n$ numbers using induction.

 $\textbf{Claim:}$ The probability that the sum is less than one is $\frac{1}{n!}$.

 $\textbf{Base Case:}$ For just $1$ number, the probability is $1$.

 $\textbf{Induction step:}$ Suppose that the probability for $n$ numbers is $\frac{1}{n!}$. We will prove that the probability for $n+1$ numbers is $\frac{1}{(n+1)!}$. To prove this, we consider that the area of an $n+1$-dimensional tetrahedron is simply the area/volume of the base times the height divided by $n+1$. 

 Of course, the area of the base is $\frac{1}{n!}$, and the height is $1$, and thus, we obtain $\frac{1}{n! \cdot (n+1)} = \frac{1}{(n+1)!}$ as our volume (this may be hard to visualize for higher dimensions). The induction step is complete.

 The probability of the sum being less than $1$ is $\frac{1}{n!}$, and the probability of the sum being more than $1$ is $\frac{n!-1}{n!}$. This trivializes the problem. The answer is 
 \[\frac{1}{2} \cdot \frac{2! - 1}{2!} + \frac{1}{2} \cdot \frac{3! - 1}{3!} = \boxed{\textbf{(C) }\frac{2}{3}}.\]




\textbf{Solution 3 (Observations)}

There are two cases: Amelia takes two steps or three steps. 

 The former case has a probability of $\frac{1}{2}$, as stated in Solution 1, and thus the latter also has a probability of $\frac{1}{2}$. 

 The probability that Amelia passes $1$ after two steps is also $\frac{1}{2}$, as it is symmetric to the probability above.

 Thus, if the probability that Amelia passes $1$ after three steps is $x$, our total probability is $\frac{1}{2} \cdot \frac{1}{2} + \frac{1}{2} \cdot x$. We know that $0 < x < 1$, and it is relatively obvious that $x > \frac{1}{2}$ (because the probability that $x > \frac{3}{2}$ is $\frac{1}{2}$). This means that our total probability is between $\frac{1}{2}$ and $\frac{3}{4}$, non-inclusive, so the only answer choice that fits is $\boxed{\textbf{(C) }\frac{2}{3}}$.

 \textbf{Remark (Calculus)} It is not immediately clear why three random numbers between $0$ and $1$ have a probability of $\frac{5}{6}$ of summing to more than $1$. Here is a proof:

 Let us start by finding the probability that two random numbers between $0$ and $1$ have a sum of more than $x$, where $0 \leq x \leq 1$. 

 Suppose that our two numbers are $y$ and $z$. Then, the probability that $y > x$ (which means that $y + z > x$) is $1 - x$, and the probability that $y < x$ is $x$.

 If $y < x$, the probability that $y + z > x$ is $1 - x + y$. This is because the probability that $y + z < x$ is equal to the probability that $z < x - y$, which is $x - y$, so our total probability is $1 - (x - y) = 1 - x + y$.

 Let us now find the average of the probability that $y + z > x$ when $y < x$. Since $y$ is a random number between $0$ and $x$, its average is $\frac{x}{2}$. Thus, our average is $1 - x + \frac{1}{2} = 1 - \frac{x}{2}$. 

 Hence, our total probability is equal to 
 \[1(1-x) + \left(1 - \frac{x}{2}\right)(x) = 1 - \frac{1}{2}x^2.\]Now, let us find the probability that three numbers uniformly distributed between $0$ and $1$ sum to more than $1$.

 Let our three numbers be $a$, $b$, and $c$. Then, the probability that $a + b + c > 1$ is equal to the probability that $b + c$ is greater than $1 - a$, which is equal to $1 - \frac{1}{2}(1 - a)^2$. 

 To find the total probability, we must average over all values of $a$. This average is simply equal to the area under the curve $1 - \frac{1}{2}(1-x)^2$ from $0$ to $1$, all divided by $1$. We can compute this value using integrals:
 \begin{align*} \frac{\int_0^1 \! 1 - \frac{1}{2}(1 - x)^2 \mathrm{d}x}{1} &= \int_0^1 \! 1 - \frac{1}{2}(1 - x)^2 \mathrm{d}x \\ &= 1 - \frac{1}{2}\int_0^1 \! (1 - x)^2 \mathrm{d}x \\ &= 1 - \frac{1}{2}\int_0^1 \! x^2 \mathrm{d}x \\ &= 1 - \frac{1}{2}\left(\frac{1}{3}\right) \\ &= \frac{5}{6}. \end{align*}
For those who don't know calculus, $\int_m^n \! f(x) \mathrm{d}x$ is the area under the curve $f(x)$ from $m$ to $n$.

 \textbf{Remark (Rigorous Calculus)} In the language of probability, we have three random variables $X$, $Y$, $Z$, each independent and uniformly distributed over the interval $(0, 1)$. We are interested in the probability that $X+Y+Z>1$, or $P(X+Y+Z>1)$.

 It follows from the definition of the joint probability density function $f(x, y, z)$ that (we assume that $0<x, y, z<1$):\begin{align}P(X+Y+Z>1)=\iiint_{x+y+z>1} f(x,y,z) \,dx\,dy\,dz.\end{align}Since $X, Y, Z$ are independent, we have $f(x, y, z)=f_X(x)*f_Y(y)*f_Z(z)$, where $f_X(x), f_Y(y), f_Z(z)$ represent the probability densities of $X, Y, Z$ respectively. Recall that $X, Y, Z$ are uniformly distributed over the interval $(0, 1)$. Hence, $f_X(x)=f_Y(y)=f_Z(z)=1$.

 It remains to evaluate the following triple integral\begin{align}\iiint_{x+y+z>1} \,dx\,dy\,dz=\frac{5}{6}.\end{align}Therefore, $P(X+Y+Z>1)=\frac{5}{6}$.

 respect to whoever would do this during the comp

 -guy\_who\_lieks\_frogs

 


	\item (\textit{AMC 12 12B 2022}) Let $x_0,x_1,x_2,\dotsc$ be a sequence of numbers, where each $x_k$ is either $0$ or $1$. For each positive integer $n$, define 

 \[S_n = \sum_{k=0}^{n-1} x_k 2^k\]
Suppose $7S_n \equiv 1 \pmod{2^n}$ for all $n \geq 1$. What is the value of the sum

 \[x_{2019} + 2x_{2020} + 4x_{2021} + 8x_{2022}?\]
$\textbf{(A) } 6 \qquad \textbf{(B) } 7 \qquad \textbf{(C) }12\qquad \textbf{(D) } 14\qquad \textbf{(E) }15$

 \textbf{Answer: }A



\textbf{Solution 1}

In binary 

 


	\item (\textit{AMC 12 12B 2022}) The figure below depicts a regular $7$-gon inscribed in a unit circle.

                    \begin{figure}[H]
                        \centering
                        \includegraphics[width=0.5\textwidth]{Problem/AMC_12_12B_2022_24_Problem_Figure_1.png}
                    \end{figure}
                    
What is the sum of the $4$th powers of the lengths of all $21$ of its edges and diagonals?


 $\textbf{(A) }49 \qquad \textbf{(B) }98 \qquad \textbf{(C) }147 \qquad \textbf{(D) }168 \qquad \textbf{(E) }196$

 \textbf{Answer: }C



\textbf{Solution 1 (Complex Numbers)}

There are $7$ segments whose lengths are $2 \sin \frac{\pi}{7}$, $7$ segments whose lengths are $2 \sin \frac{2 \pi}{7}$, $7$ segments whose lengths are $2 \sin \frac{3\pi}{7}$.

 Therefore, the sum of the $4$th powers of these lengths is
 \begin{align*} 7 \cdot 2^4 \sin^4 \frac{\pi}{7} + 7 \cdot 2^4 \sin^4 \frac{2 \pi}{7} + 7 \cdot 2^4 \sin^4 \frac{3 \pi}{7} & = \frac{7 \cdot 2^4}{(2i)^4} \left( e^{i \frac{\pi}{7}} - e^{i \frac{-\pi}{7}} \right)^4 + \frac{7 \cdot 2^4}{(2i)^4} \left( e^{i \frac{2 \pi}{7}} - e^{i \frac{-2 \pi}{7}} \right)^4 + \frac{7 \cdot 2^4}{(2i)^4} \left( e^{i \frac{3 \pi}{7}} - e^{i \frac{-3 \pi}{7}} \right)^4 \\ & = 7 \left( e^{i \frac{4 \pi}{7}} - 4 e^{i \frac{2 \pi}{7}} + 6 - 4 e^{- i \frac{2 \pi}{7}} + e^{- i \frac{4 \pi}{7}} \right) \\ & \quad + 7 \left( e^{i \frac{8 \pi}{7}} - 4 e^{i \frac{4 \pi}{7}} + 6 - 4 e^{- i \frac{4 \pi}{7}} + e^{- i \frac{8 \pi}{7}} \right) \\ & \quad + 7 \left( e^{i \frac{12 \pi}{7}} - 4 e^{i \frac{6 \pi}{7}} + 6 - 4 e^{- i \frac{6 \pi}{7}} + e^{- i \frac{12 \pi}{7}} \right) \\ & = 7 \left( e^{i \frac{4 \pi}{7}} + e^{i \frac{8 \pi}{7}} + e^{i \frac{12 \pi}{7}} + e^{-i \frac{4 \pi}{7}} + e^{-i \frac{8 \pi}{7}} + e^{-i \frac{12 \pi}{7}} \right) \\ & \quad - 7 \cdot 4 \left( e^{i \frac{2 \pi}{7}} + e^{i \frac{4 \pi}{7}} + e^{i \frac{6 \pi}{7}} + e^{-i \frac{2 \pi}{7}} + e^{-i \frac{4 \pi}{7}} + e^{-i \frac{6 \pi}{7}} \right) \\ & \quad + 7 \cdot 6 \cdot 3 \\ & = 7 \left( e^{i \frac{4 \pi}{7}} + e^{-i \frac{6 \pi}{7}} + e^{-i \frac{2 \pi}{7}} + e^{-i \frac{4 \pi}{7}} + e^{i \frac{6 \pi}{7}} + e^{i \frac{2 \pi}{7}} \right) \\ & \quad - 7 \cdot 4 \left( e^{i \frac{2 \pi}{7}} + e^{i \frac{4 \pi}{7}} + e^{i \frac{6 \pi}{7}} + e^{-i \frac{2 \pi}{7}} + e^{-i \frac{4 \pi}{7}} + e^{-i \frac{6 \pi}{7}} \right) \\ & \quad + 7 \cdot 6 \cdot 3 \\ & = -7 + 7 \cdot 4 + 7 \cdot 6 \cdot 3 \\ & = \boxed{\textbf{(C) }147}, \end{align*}
where the fourth from the last equality follows from the property that
 \begin{align*} e^{i \frac{2 \pi}{7}} + e^{i \frac{4 \pi}{7}} + e^{i \frac{6 \pi}{7}} + e^{-i \frac{2 \pi}{7}} + e^{-i \frac{4 \pi}{7}} + e^{-i \frac{6 \pi}{7}} & = e^{-i \frac{6 \pi}{7}} \sum_{j=0}^6 e^{i \frac{2 \pi j}{7}} - 1  \\ & = 0 - 1 \\ & = -1 . \end{align*}
~Steven Chen (Professor Chen Education Palace, www.professorchenedu.com)

 



\textbf{Solution 2 (Trigonometry)}

There are $7$ segments whose lengths are $2 \sin \frac{\pi}{7}$, $7$ segments whose lengths are $2 \sin \frac{2 \pi}{7}$, $7$ segments whose lengths are $2 \sin \frac{3\pi}{7}$.

 Therefore, the sum of the $4$th powers of these lengths is
 \begin{align*} & 7 \cdot 2^4 \sin^4 \frac{\pi}{7} + 7 \cdot 2^4 \sin^4 \frac{2 \pi}{7} + 7 \cdot 2^4 \sin^4 \frac{3 \pi}{7} \\ & = 7 \cdot 2^4 \left( \frac{1 - \cos \frac{2 \pi}{7}}{2} \right)^2 + 7 \cdot 2^4 \left( \frac{1 - \cos \frac{4 \pi}{7}}{2} \right)^2 + 7 \cdot 2^4 \left( \frac{1 - \cos \frac{6 \pi}{7}}{2} \right)^2 \\ & = 7 \cdot 2^2 \left( 1 - 2 \cos \frac{2 \pi}{7} + \cos^2 \frac{2 \pi}{7} \right) + 7 \cdot 2^2 \left( 1 - 2 \cos \frac{4 \pi}{7} + \cos^2 \frac{4 \pi}{7} \right) + 7 \cdot 2^2 \left( 1 - 2 \cos \frac{6 \pi}{7} + \cos^2 \frac{6 \pi}{7} \right) \\ & = 7 \cdot 2^2 \cdot 3 - 7 \cdot 2^3 \left( \cos \frac{2 \pi}{7} + \cos \frac{4 \pi}{7} + \cos \frac{6 \pi}{7} \right) + 7 \cdot 2^2 \left( \cos^2 \frac{2 \pi}{7} + \cos^2 \frac{4 \pi}{7}  + \cos^2 \frac{6 \pi}{7} \right) \\ & = 7 \cdot 2^2 \cdot 3 - 7 \cdot 2^3 \left( \cos \frac{2 \pi}{7} + \cos \frac{4 \pi}{7} + \cos \frac{6 \pi}{7} \right) + 7 \cdot 2^2 \left( \frac{1 + \cos \frac{4 \pi}{7} }{2} + \frac{1 + \cos \frac{8 \pi}{7} }{2} + \frac{1 + \cos \frac{12 \pi}{7} }{2} \right) \\ & = 7 \cdot 2^2 \cdot 3 - 7 \cdot 2^3 \left( \cos \frac{2 \pi}{7} + \cos \frac{4 \pi}{7} + \cos \frac{6 \pi}{7} \right) + 7 \cdot 2 \cdot 3 + 7 \cdot 2 \left( \cos \frac{4 \pi}{7} + \cos \frac{8 \pi}{7} + \cos \frac{12 \pi}{7} \right) \\ & = 7 \cdot 2^2 \cdot 3 - 7 \cdot 2^3 \left( \cos \frac{2 \pi}{7} + \cos \frac{4 \pi}{7} + \cos \frac{6 \pi}{7} \right) + 7 \cdot 2 \cdot 3 + 7 \cdot 2 \left( \cos \frac{4 \pi}{7} + \cos \frac{6 \pi}{7} + \cos \frac{2 \pi}{7} \right) \\ & = 7 \cdot 2 \cdot 3 \left( 2 + 1 \right) - 7 \cdot 2 \left( 4 - 1 \right) \left( \cos \frac{4 \pi}{7} + \cos \frac{6 \pi}{7} + \cos \frac{2 \pi}{7} \right) \\ & = 7 \cdot 2 \cdot 3 \left( 2 + 1 \right) - 7 \cdot 2 \left( 4 - 1 \right) \cdot \left( - \frac{1}{2} \right) \\ & = \boxed{\textbf{(C) }147}, \end{align*}
where the second from the last equality follows from the property that
 \[ \cos \frac{4 \pi}{7} + \cos \frac{6 \pi}{7} + \cos \frac{2 \pi}{7} = - \frac{1}{2} . \]

 For explanation see \href{https://artofproblemsolving.com/wiki/index.php/2022\_AMC\_12B\_Problems/Problem\_24\#Supplement\_.28Explanation\_to\_Why\_cos.282.CF.80.2F7.29\_.2B\_cos.284.CF.80.2F7.29\_.2B\_cos.286.CF.80.2F7.29\_.3D\_-1.2F2.29}{supplement}.

 



\textbf{Solution 3 (Complex Numbers and Trigonometry)}

As explained in Solutions 1 and 2, what we are trying to find is $7 \cdot 2^4 \sin^4 \frac{\pi}{7} + 7 \cdot 2^4 \sin^4 \frac{2 \pi}{7} + 7 \cdot 2^4 \sin^4 \frac{3 \pi}{7}$. Using trig we get 
 \begin{align*}  & \sin^4 \frac{\pi}{7} + \sin^4 \frac{2 \pi}{7} + \sin^4 \frac{3 \pi}{7} \\ = & \sin^2 \frac{\pi}{7} \left(1 - \cos^2 \frac{\pi}{7} \right) + \sin^2 \frac{2\pi}{7} \left(1 - \cos^2 \frac{2\pi}{7} \right) + \sin^2 \frac{3\pi}{7} \left(1 - \cos^2 \frac{3\pi}{7} \right) \\ = & \sin^2 \frac{\pi}{7} - \left(\frac{1}{2} \sin \frac{2\pi}{7}\right)^2 + \sin^2 \frac{2\pi}{7} - \left(\frac{1}{2} \sin \frac{4\pi}{7}\right)^2 + \sin^2 \frac{3\pi}{7} - \left(\frac{1}{2} \sin \frac{6\pi}{7}\right)^2\\ = & \sin^2 \frac{\pi}{7} - \frac{1}{4} \sin^2 \frac{2\pi}{7} + \sin^2 \frac{2\pi}{7} - \frac{1}{4} \sin^2 \frac{4\pi}{7} + \sin^2 \frac{3\pi}{7} - \frac{1}{4} \sin^2 \frac{6\pi}{7} \\ = & \frac{3}{4} \left(\sin^2 \frac{\pi}{7} + \sin^2 \frac{2\pi}{7} + \sin^2 \frac{3\pi}{7}\right) \\ = & \frac{3}{4} \cdot \frac{1}{2} \left(1 - \cos \frac{2\pi}{7} + 1 - \cos \frac{4\pi}{7} + 1 - \cos \frac{6\pi}{7} \right)\\ = & \frac{3}{4} \cdot \frac{1}{2} \left(3 - \left(-\frac{1}{2}\right)\right) \\ = & \frac{21}{16}. \end{align*}
Like in the second solution, we also use the fact that $\cos \frac{2\pi}{7} + \cos \frac{4\pi}{7} + \cos \frac{6\pi}{7} = -\frac{1}{2}$, which admittedly might need some explanation. 

 For explanation see \href{https://artofproblemsolving.com/wiki/index.php/2022\_AMC\_12B\_Problems/Problem\_24\#Supplement\_.28Explanation\_to\_Why\_cos.282.CF.80.2F7.29\_.2B\_cos.284.CF.80.2F7.29\_.2B\_cos.286.CF.80.2F7.29\_.3D\_-1.2F2.29}{supplement}.

 Notice that 
 \begin{align*} \cos \frac{2\pi}{7} + \cos \frac{4\pi}{7} + \cos \frac{6\pi}{7} & = \frac{1}{2}\left(e^\frac{2i\pi}{7}+ e^{-\frac{2i\pi}{7}} \right) + \frac{1}{2}\left( e^\frac{4i\pi}{7}+ e^{-\frac{4i\pi}{7}} \right) + \frac{1}{2}\left( e^\frac{6i\pi}{7}+ e^{-\frac{6i\pi}{7}}\right) \\ & = \frac{1}{2}\left(e^\frac{2i\pi}{7}+ e^{-\frac{2i\pi}{7}}+ e^\frac{4i\pi}{7}+ e^{-\frac{4i\pi}{7}} +e^\frac{6i\pi}{7}+ e^{-\frac{6i\pi}{7}} + 1\right) - \frac{1}{2} \end{align*}
In the brackets we have the sum of the roots of the polynomial $x^7 - 1 = 0$. These sum to $0$ by \href{https://artofproblemsolving.com/wiki/index.php?title=Vieta\%E2\%80\%99s\_formulas&action=edit&redlink=1}{Vieta’s formulas}, and the desired identity follows. See \href{https://artofproblemsolving.com/wiki/index.php?title=Roots\_of\_unity}{Roots of unity} if you have not seen this technique.

 Going back to the question: 
 \[7 \cdot 2^4 \sin^4 \frac{\pi}{7} + 7 \cdot 2^4 \sin^4 \frac{2 \pi}{7} + 7 \cdot 2^4 \sin^4 \frac{3 \pi}{7} = 7 \cdot 2^4 \left(\sin^4 \frac{\pi}{7} + \sin^4 \frac{2 \pi}{7} + \sin^4 \frac{3 \pi}{7}\right) = 7 \cdot 2^4 \cdot \frac{21}{16} = \boxed{\textbf{(C) }147}.\]~obscene\_kangaroo

 



\textbf{Solution 4 (Trigonometry)}

This solution follows the same steps as the trigonometry solutions (Solutions 2 and 3), except it gives an alternate way to prove the statement below true without complex numbers:

 
 \[\cos \frac{2\pi}{7} + \cos \frac{4\pi}{7} + \cos \frac{6\pi}{7} = -\frac{1}{2}\]

 \begin{align*} S &= \cos \frac{2\pi}{7} + \cos \frac{4\pi}{7} + \cos \frac{6\pi}{7}, \\ S^2 &= \cos^2 \frac{2\pi}{7} + \cos^2 \frac{4\pi}{7} + \cos^2 \frac{6\pi}{7} + 2\cos \frac{2\pi}{7}\cos \frac{4\pi}{7} + 2\cos \frac{2\pi}{7}\cos \frac{6\pi}{7} + 2\cos \frac{4\pi}{7}\cos \frac{6\pi}{7} \\ &= \left(\frac{1+ \cos \frac{4\pi}{7}}{2}\right) + \left(\frac{1+ \cos \frac{8\pi}{7}}{2}\right) + \left(\frac{1+ \cos \frac{12\pi}{7}}{2}\right) + 2\cos \frac{2\pi}{7}\cos \frac{4\pi}{7} + 2\cos \frac{2\pi}{7}\cos \frac{6\pi}{7} + 2\cos \frac{4\pi}{7}\cos \frac{6\pi}{7} \\ &= \frac{1}{2}(3 + S)  + \left(\cos \frac{6\pi}{7} + \cos \frac{2\pi}{7}\right) + \left(\cos \frac{8\pi}{7} + \cos \frac{4\pi}{7}\right) + \left(\cos \frac{10\pi}{7} + \cos \frac{2\pi}{7}\right) \\ &= \frac{1}{2}(3 + S) + 2\cos \frac{2\pi}{7} + 2\cos \frac{4\pi}{7} + 2\cos \frac{6\pi}{7}\\ &= \frac{1}{2}(3 + S) + 2S. \\ \end{align*}
We end up with 
 \[2S^2 - 5S - 3 = 0.\]Using the quadratic formula, we find the solutions for $S$ to be $-\frac{1}{2}$ and $3$. Because $3$ is impossible, $S = -\frac{1}{2}$.  With this result, following similar to steps to Solutions 2 and 3 will get $\boxed{\textbf{(C) }147}$




\textbf{Solution 5 (Law of Cosines)}

Let $x$, $y$, and $z$ be the lengths of the chords with arcs $\frac{2\pi}{7}$, $\frac{4\pi}{7}$ and $\frac{6\pi}{7}$ respectively. 

 Then by the law of cosines we get:
 \begin{align*} x^2 &= 2\left(1-\cos\frac{2\pi}{7}\right), \\ y^2 &= 2\left(1-\cos\frac{4\pi}{7}\right), \\ z^2 &= 2\left(1-\cos\frac{6\pi}{7}\right). \end{align*}
The answer is then just $7(x^4+y^4+z^4)$ (since there's $7$ of each diagonal/side), obtained by summing the squares of the above equations and then multiplying by $7$. 
 \begin{align*} & \quad7\cdot2^2\left( \left(1-\cos\frac{2\pi}{7}\right)^2 + \left(1-\cos\frac{4\pi}{7}\right)^2 + \left(1-\cos\frac{6\pi}{7}\right)^2 \right) \\ &= 7\cdot4\left( \left(1-2\cos\frac{2\pi}{7}+\cos^2\frac{2\pi}{7}\right) + \left(1-2\cos\frac{4\pi}{7}+\cos^2\frac{4\pi}{7}\right) + \left(1-2\cos\frac{6\pi}{7}+\cos^2\frac{6\pi}{7}\right) \right) \\ &= 7\cdot4\left(3-2\left(\cos\frac{2\pi}{7}+\cos\frac{4\pi}{7}+\cos\frac{6\pi}{7}\right) + \left(\cos^2\frac{2\pi}{7}+\cos^2\frac{4\pi}{7}+\cos^2\frac{6\pi}{7}\right) \right) \\ &= 7\cdot4\left(3-2\left(\cos\frac{2\pi}{7}+\cos\frac{4\pi}{7}+\cos\frac{6\pi}{7}\right) + \frac12 \left(1+\cos\frac{4\pi}{7}+1+\cos\frac{8\pi}{7}+1+\cos\frac{12\pi}{7}\right) \right) \\ &= 7\cdot4\left(\frac{9}{2}-2\left(\cos\frac{2\pi}{7}+\cos\frac{4\pi}{7}+\cos\frac{6\pi}{7}\right) + \frac12 \left(\cos\frac{4\pi}{7}+\cos\frac{6\pi}{7}+\cos\frac{2\pi}{7}\right) \right) \\ &= 7\cdot4\left(\frac{9}{2}-\frac{3}{2}\left(\cos\frac{2\pi}{7}+\cos\frac{4\pi}{7}+\cos\frac{6\pi}{7}\right) \right) \\ &= 7\cdot4\left(\frac{9}{2}-\frac{3}{2}\left(-\frac{1}{2}\right)\right) \\ &= \boxed{\textbf{(C) }147}. \end{align*}


 (Use the identity that $\cos\frac{2\pi}{7}+\cos\frac{4\pi}{7}+\cos\frac{6\pi}{7} = -\frac{1}{2}$.)

 For explanation of this identity see \href{https://artofproblemsolving.com/wiki/index.php/2022\_AMC\_12B\_Problems/Problem\_24\#Supplement\_.28Explanation\_to\_Why\_cos.282.CF.80.2F7.29\_.2B\_cos.284.CF.80.2F7.29\_.2B\_cos.286.CF.80.2F7.29\_.3D\_-1.2F2.29}{supplement}.

 - SAHANWIJETUNGA

 



\textbf{Solution 6 (Ruler Measure)}

Hope you had a ruler handy! This problem can be done with a ruler and basic estimation.

 First, measuring the radius of the circle obtains $2.9$ cm (when done on the paper version). Thus, any other measurement we get for the sides/diagonals should be divided by $2.9$.

 Measuring the sides of the circle gets $2.5$ cm. The shorter diagonals are $4.5$ cm, and the longest diagonals measure $5.6$ cm. Thus, we'd like to estimate 
 \[7\left(\frac{2.5}{2.9}\right)^4 + 7\left(\frac{4.5}{2.9}\right)^4 + 7\left(\frac{5.6}{2.9}\right)^4.\]

 We know $\left(\frac{2.5}{2.9}\right)^4$ is slightly less than $1.$ Let's approximate it as 1 for now. Thus, $7\left(\frac{2.5}{2.9}\right)^4 \approx 7.$

 Next, $\left(\frac{4.5}{2.9}\right)^4$ is slightly more than $\left(\frac{4.5}{3}\right)^4.$ We know $\left(\frac{4.5}{3}\right)^4 = 1.5^4 = \frac{81}{16},$ slightly more than $5,$ so we can approximate $\left(\frac{4.5}{2.9}\right)^4$ as $5.5.$ Thus, $7\left(\frac{4.5}{2.9}\right)^4 \approx 38.5.$

 Finally, $\left(\frac{5.6}{2.9}\right)^4$ is slightly less than $\left(\frac{5.6}{2.8}\right)^4 = 2^4 = 16.$ We say it's around $15,$ so then $7\left(\frac{5.6}{2.9}\right)^4 \approx 105.$

 Adding what we have, we get $105 + 38.5 + 1 = 144.5$ as our estimate. We see $\boxed{\textbf{(C) }147}$ is very close to our estimate, so we have successfully finished the problem.

 



\textbf{Solution 7 (Pythagorean Theorem and Trig)}

First, we put the figure in the coordinate plane with the center of the circle at the origin and a vertex on the positive x-axis. Thus, the coordinates of the vertices will be the terminal points of integer multiples of the angle $\frac{2\pi}{7},$ which are 
 \[\left(\cos\dfrac{2\pi n}{7}, \sin\dfrac{2\pi n}{7}\right)\] for integer $n.$ Then, we notice there are three types of diagonals: the ones with chords of arcs $\dfrac{2\pi}{7}, \dfrac{4\pi}{7},$ and $\dfrac{6\pi}{7}.$ We notive there are here are $7$ of each type of diagonal. Then, we use the \href{https://artofproblemsolving.com/wiki/index.php?title=Pythagorean\_theorem}{pythagorean theorem} to find the distance from $\left(\cos\frac{2\pi n}{7}, \sin\frac{2\pi n}{7}\right)$ to $(1, 0)$:
 \[\left(\sqrt{\left(\cos\frac{2\pi n}{7}-1\right)^2+\sin^2\frac{2\pi n}{7}}\right)^4=\left(\sqrt{\cos^2\frac{2\pi n}{7}+\sin^2\frac{2\pi n}{7}-2\cos\frac{2\pi n}{7}+1}\right)^4\]
 \[=\left(\sqrt{2-2\cos\frac{2\pi n}{7}}\right)^4\]
 \[=4\cos^2\frac{2\pi n}{7}-8\cos\frac{2\pi n}{7}+4.\]By the cosine double angle identity, $\cos{2\theta}=2\cos^2\theta-1.$ This means that $2\cos^2\theta=\cos{2\theta}+1.$ Substituting this in,
 \[4\cos^2\frac{2\pi n}{7}-8\cos\frac{2\pi n}{7}+4=2\left(\cos\frac{4\pi n}{7}+1\right)-8\cos\frac{2\pi n}{7}+4=2\cos\frac{4\pi n}{7}-8\cos\frac{2\pi n}{7}+6.\] Summing this up for $n=1,2,3,$ 
 \[2\cos\frac{4\pi}{7}-8\cos\frac{2\pi}{7}+6+2\cos\frac{8\pi}{7}-8\cos\frac{4\pi}{7}+6+2\cos\frac{12\pi}{7}-8\cos\frac{6\pi}{7}+6\] 
 \[=2\left(\cos\frac{4\pi}{7}+\cos\frac{8\pi}{7}+\cos\frac{12\pi}{7}\right)-8\left(\cos\frac{2\pi}{7}+\cos\frac{4\pi}{7}+\cos\frac{6\pi}{7}\right)+18\]
 \[=2\left(\cos\frac{4\pi}{7}+\cos\frac{6\pi}{7}+\cos\frac{2\pi}{7}\right)-8\left(\cos\frac{2\pi}{7}+\cos\frac{4\pi}{7}+\cos\frac{6\pi}{7}\right)+18\]
 \[=2(-0.5)-8(-0.5)+18=21.\](These equalities are based on $\cos\theta=\cos(2\pi-\theta)$ and $\cos\frac{2\pi}{7}+\cos\frac{4\pi}{7}+\cos\frac{6\pi}{7}=-\frac{1}{2}.$) (For explanation of this see \href{https://artofproblemsolving.com/wiki/index.php/2022\_AMC\_12B\_Problems/Problem\_24\#Supplement\_.28Explanation\_to\_Why\_cos.282.CF.80.2F7.29\_.2B\_cos.284.CF.80.2F7.29\_.2B\_cos.286.CF.80.2F7.29\_.3D\_-1.2F2.29}{supplement}.)Finally, because there are $7$ of each type of diagonal, the answer is $7\cdot 21=\boxed{147}.$




\textbf{Solution 8 (Roots of Unity)}

Place the figure on the complex plane and let $w = e^{\frac{2\pi}{7}}$ (a $7$th root of unity). The vertices of the $7$-gon are $w^0,w^1,w^2,\dots,w^6$. We wish to find
 \[\sum_{i=0}^5\sum_{j=i+1}^6\lvert w^i-w^j\rvert^4 = \frac{1}{2}\sum_{i=0}^6\sum_{j=0}^6\lvert w^i-w^j\rvert^4.\]The second expression is more convenient to work with. The factor of $\frac{1}{2}$ is because it double-counts each edge (and includes when $i=j$, but the length is $0$ so it doesn't matter).

 Recall the identity $|z|^2 = z\overline{z}$. Since $|w| = 1$, $\overline{w} = w^{-1}$, and
 \begin{align*}     \lvert w^i-w^j\rvert^4 &= ((w^i-w^j)(w^{-i}-w^{-j}))^2 \\     &= (2-w^{i-j}-w^{j-i})^2 \\     &= 4+w^{2i-2j}+w^{2j-2i}+2(-2w^{i-j}-2w^{j-i}+1) \\     &= 6-4w^{i-j}-4w^{j-i}+w^{2i-2j}+w^{2j-2i}. \end{align*}


 But notice that, for a fixed value of $i$, over all values of $j$ from $0$ to $6$, by properties of modular arithmetic with $w^7 = w^0$ (or expanding),each of the $w^\bullet$-terms takes on every value in $w^0,w^1,w^2,\dots,w^6$ exactly once. Since $\sum_{j=0}^6 w^j = 0$,
 \begin{align*}     \frac{1}{2}\sum_{i=0}^6\sum_{j=0}^6\lvert w^i-w^j\rvert^4 &= \frac{1}{2}\sum_{i=0}^6\sum_{j=0}^6(6-4w^{i-j}-4w^{j-i}+w^{2i-2j}+w^{2j-2i}) \\     &= \frac{1}{2}\sum_{i=0}^6\left(\sum_{j=0}^6 6-4\sum_{j=0}^6 w^{i-j}-4\sum_{j=0}^6 w^{j-i}+\sum_{j=0}^6 w^{2i-2j}+\sum_{j=0}^6 w^{2j-2i}\right) \\     &= \frac{1}{2}\sum_{i=0}^6\left(\sum_{j=0}^6 6-4\sum_{j=0}^6 w^j-4\sum_{j=0}^6 w^j+\sum_{j=0}^6 w^j+\sum_{j=0}^6 w^j\right) \\     &= \frac{1}{2}\sum_{i=0}^6\sum_{j=0}^6 6 \\     &= \frac{1}{2}\cdot7\cdot7\cdot6 \\     &= \boxed{\textbf{(C) }147}. \end{align*}


 Fun fact: This generalizes to any regular $n$-gon with $n \ge 3$ to obtain
 \[\frac{1}{2}\sum_{i=0}^{n-1}\sum_{j=0}^{n-1}6 = 3n^2,\]though for even $n$, $w^{2i-2j}$ and $w^{2j-2i}$ take on every value in $w^0,w^2,w^4,\dots,w^{n-2}$ exactly twice. These are the $\frac{n}{2}$th roots of unity, so they still sum to $0$.

 This also explains why $n < 3$ doesn't follow the pattern. For $n = 2$ (one edge of length $2$), since $w^{2i-2j} = (w^2)^{i-j} = 1^{i-j} = 1$ and similarly for $w^{2j-2i}$, the constant term in the summation becomes $8$ instead of $6$. For $n = 1$ (no edges), the constant term becomes $0$.

 -maxamc

 



\textbf{Solution 9 (Inductive Reasoning)}

This is how I solve this problem:

 It's easy to solve for $3$-gon, $4$-gon, and $6$-gon inscribed in a unit circle. (Okay, it's just the weird names for triangle, square, and hexagon)

 For $3$-gon, the sum is equal to $3$ times the $4$th power of an edge. Thus, 
 \[S_3=3\,\cdot\,\left(\sqrt{3}\right)^4=27.\]

 For $4$-gon, the sum is equal to $4$ times the $4$th power of an edge, and $2$ times the $4$th power of the diagonal. Thus, 
 \[S_4=4\,\cdot\,\left(\sqrt{2}\right)^4+2\,\cdot\,\left(2\right)^4=48.\]

 For $6$-gon, the sum is equal to $6$ times the $4$th power of an edge, $6$ times the $4$th power of the short diagonal, and $3$ times the $4$th power of the long diagonal. Thus, 
 \[S_6=6\,\cdot\,\left(1\right)^4+6\,\cdot\,\left(\sqrt{3}\right)^4+3\,\cdot\,\left(2\right)^4=108.\]

 Then, I quickly noticed that $27=3\,\cdot\,3^2$, $48=3\,\cdot\,4^2$, and $108=3\,\cdot\,6^2$. So reasonably, it will work out this formula, $S_n=3n^2$. (This step is purely out of guessing, maybe have a look at Solution 8 for more info...)

 By inductive reasoning, we got $S_7=3\,\cdot\,7^2=\boxed{\textbf{(C) }147}$.

 - Prof. Joker

 



\textbf{Solution 10 (cheese)}

There are $7$ segments with length $2\sin{\frac{\pi}{7}}$, $7$ segments with length $2\sin{\frac{2\pi}{7}}$, and $7$ segments with length $2\sin{\frac{3\pi}{7}}$. The sum is $112(\sin^4\frac{\pi}{7}+\sin^4\frac{2\pi}{7}+\sin^4\frac{3\pi}{7})$ Note that $\frac{3}{7} = 0.425871...$, so $\sin\frac{3\pi}{7} \approx \sin\frac{\pi}{2}$, but a little less (The actual value is approximately $0.975$), and based on the answer choices, we can conclude that our answer will be a little greater than $1$, so we make the educated guess of $\boxed{\textbf{(C) }147}$

 \textbf{Supplement (Explanation of why cos(2$\pi$/7) + cos(4$\pi$/7) + cos(6$\pi$/7) = -1/2)} The sum of all roots of unity is $0$, therefore, the sum of the real parts of the roots of unity is $0$. 

 Consider the roots of unity of $x^n-1=0$ where $n$ is odd on the unit circle, the cosine (real parts) of all the roots under the x-axis is the same as the cosine of all the roots above the x-axis because of symmetry. Hence, $2$ times the sum of the cosines of all the roots above the x-axis plus $1$ is $0$.

 In the context of this problem, $2(\cos\frac{2\pi}{7}+\cos\frac{4\pi}{7}+\cos\frac{6\pi}{7}) + 1 = 0$. Hence, $\cos\frac{2\pi}{7}+\cos\frac{4\pi}{7}+\cos\frac{6\pi}{7} = -\frac{1}{2}$



	\item (\textit{AMC 12 12B 2022}) Four regular hexagons surround a square with a side length $1$, each one sharing an edge with the square, as shown in the figure below. The area of the resulting 12-sided outer nonconvex polygon can be written as $m\sqrt{n} + p$, where $m$, $n$, and $p$ are integers and $n$ is not divisible by the square of any prime. What is $m + n + p$?


 
                    \begin{figure}[H]
                        \centering
                        \includegraphics[width=0.5\textwidth]{Problem/AMC_12_12B_2022_25_Problem_Figure_1.png}
                    \end{figure}
                    
$\textbf{(A) } -12 \qquad \textbf{(B) }-4 \qquad  \textbf{(C) } 4 \qquad \textbf{(D) }24 \qquad \textbf{(E) }32$

 \textbf{Answer: }B



\textbf{Solution 1 (Coord bash)}


                    \begin{figure}[H]
                        \centering
                        \includegraphics[width=0.5\textwidth]{Solution/AMC_12_12B_2022_25_Solution_Figure_1.png}
                    \end{figure}
                    

 Refer to the diagram above.

 Let the origin be at the center of the square, $A$ be the intersection of the top and right hexagons, $B$ be the intersection of the top and left hexagons, and $M$ and $N$ be the top points in the diagram.

 By symmetry, $A$ lies on the line $y = x$. The equation of line $AN$ is $y = -x\sqrt{3} + \frac{3}{2}\sqrt{3} - \frac{1}{2}$ (due to it being one of the sides of the top hexagon). Thus, we can solve for the coordinates of $A$ by finding the intersection of the two lines:
 \[x = -x\sqrt{3} + \frac{3\sqrt{3} - 1}{2}\]
 \[x(\sqrt{3} + 1) = \frac{3\sqrt{3} - 1}{2}\]
 \[x = \frac{3\sqrt{3}-1}{2} \cdot \frac{1}{\sqrt{3} + 1}\]
 \[= \frac{3\sqrt{3}-1}{2(\sqrt{3} + 1)} \cdot \frac{\sqrt{3} - 1}{\sqrt{3} - 1}\]
 \[= \frac{10 - 4\sqrt{3}}{4}\]
 \[= \frac{5}{2} - \sqrt{3}\]
 \[\therefore A = \left(\frac{5}{2} - \sqrt{3}, \frac{5}{2} - \sqrt{3}\right).\]

 This means that we can find the length $AB$, which is equal to $2(\frac{5}{2} - \sqrt{3}) = 5 - 2\sqrt{3}$. We will next find the area of trapezoid $ABMN$. The lengths of the bases are $1$ and $5 - 2\sqrt{3}$, and the height is equal to the $y$-coordinate of $M$ minus the $y$-coordinate of $A$. The height of the hexagon is $\sqrt{3}$ and the bottom of the hexagon lies on the line $y = \frac{1}{2}$. Thus, the $y$-coordinate of $M$ is $\sqrt{3} - \frac{1}{2}$, and the height is $2\sqrt{3} - 3$. We can now find the area of the trapezoid:
 \[[ABMN] = (2\sqrt{3} - 3)\left(\frac{1 + 5 - 2\sqrt{3}}{2}\right)\]
 \[= (2\sqrt{3} - 3)(3 - \sqrt{3})\]
 \[= 6\sqrt{3} + 3\sqrt{3} - 9 - 6\]
 \[= 9\sqrt{3} - 15.\]

 The total area of the figure is the area of a square with side length $AB$ plus four times the area of this trapezoid:
 \[\textrm{Area} = (5 - 2\sqrt{3})^2 + 4(9\sqrt{3} - 15)\]
 \[= 37 - 20\sqrt{3} + 36\sqrt{3} - 60\]
 \[= 16\sqrt{3} - 23.\]

 Our answer is $16 + 3 - 23 = \boxed{\textbf{(B) }-4}$.

 



\textbf{Solution 2}


                    \begin{figure}[H]
                        \centering
                        \includegraphics[width=0.5\textwidth]{Solution/AMC_12_12B_2022_25_Solution_Figure_2.png}
                    \end{figure}
                    

 Begin by dividing the figure as shown above. Clearly, the entire figure has 8-fold symmetry. Therefore, we can calculate the area of $ODBA$ and multiply it by 8. We split $[ODBA]$ into $[ODBC]+[ABC]$.

 Knowing the side length of the hexagon is $1$, we can use 30-60-90 triangles within the hexagon to find the total distance between opposite edges is $2\cdot \frac{\sqrt{3}}{2}=\sqrt{3}.$ Thus, $OD=\sqrt{3}-\frac{1}{2}$ and $BC=\sqrt{3}-1.$ Recognizing $DB=\frac{1}{2}$ and $ODBC$ is a trapezoid, 
 \[8\cdot[ODBC]=2\left(\sqrt{3}-\frac{1}{2}+\sqrt{3}-1\right)=4\sqrt{3}-3.\]

 Next, we aim to find $[ABC]$. By angle chasing, we find $\angle A=105^\circ,$ $\angle B=30^\circ,$ and $\angle C = 45^\circ.$ We can use the law of sines to find $AB$:

 
 \[\frac{\sin(105^\circ)}{\sqrt{3}-1}=\frac{\sin(45^\circ)}{AB}\implies AB=\frac{\sqrt{6}-\sqrt{2}}{2\sin(105^\circ)}.\]

 We may not know what $\sin(105^\circ)$ is by memory, but we can cleverly calculate it using a common trig identity:

 
 \begin{align*}     \sin(105^\circ)&=\sin(60^\circ+45^\circ),\\     &=\sin(60^\circ)\cos(45^\circ)+\cos(60^\circ)\sin(45^\circ),\\     &=\frac{\sqrt{3}}{2}\cdot\frac{\sqrt{2}}{2}+\frac{1}{2}\cdot\frac{\sqrt{2}}{2},\\     &=\frac{\sqrt{6}+\sqrt{2}}{4}. \end{align*}


 With some simplification, we'll find $AB=4-2\sqrt{3}$. Now, we can easily calculate $8\cdot [ABC]$ as 
 \[8\cdot\frac{1}{2}\cdot(\sqrt{3}-1)(4-2\sqrt{3})\sin(30^\circ)=12\sqrt{3}-20.\]

 Thus, the area of the dodecagon is $8\cdot [ODBA] = 8\cdot [ODBC] + 8 \cdot [ABC] =4\sqrt{3}-3+12\sqrt{3}-20=16\sqrt{3}-23.$

 Finally, we find 
 \[16+3-23=\boxed{\textbf{(B)}\ -4}.\]




\textbf{Solution 3 (Euclidean Geometry)}


                    \begin{figure}[H]
                        \centering
                        \includegraphics[width=0.5\textwidth]{Solution/AMC_12_12B_2022_25_Solution_Figure_3.png}
                    \end{figure}
                    

 Divide this polygon into a 'cross' and $8$ triangles like $\triangle ABC$. $BD = 2 \cdot \frac{ \sqrt{3} }{2} = \sqrt{3}$, $BE = 2BD - CD = 2\sqrt{3} - 1$, $EF = 1$, $[BEFG] = 1 \cdot (2\sqrt{3} - 1) = 2\sqrt{3} - 1$, $[\text{cross}] = 2\cdot [BEFG] - 1 = 2(2\sqrt{3} - 1)- 1 = 4\sqrt{3} - 3$

                    \begin{figure}[H]
                        \centering
                        \includegraphics[width=0.5\textwidth]{Solution/AMC_12_12B_2022_25_Solution_Figure_4.png}
                    \end{figure}
                    $\angle ABC = 120^{\circ} - 90^{\circ} = 30^{\circ}$, $\angle ACB = \frac{90^{\circ}}{2} = 45^{\circ}$, $\angle BAC = 180^{\circ} - 30^{\circ} - 45^{\circ} = 105^{\circ}$, $BC = BD - CD = \sqrt{3} - 1$. Draw a perpendicular line $AG$ to side $BC$.

 $BG + CG = BC$, $AG \cdot \sqrt{3} + AG = \sqrt{3} - 1$, $AG = \frac{\sqrt{3} - 1}{\sqrt{3} + 1} = \frac{(\sqrt{3} - 1)^2}{2} = \frac{4 - 2\sqrt{3}}{2} = 2 - \sqrt{3}$, $[ABC] = \frac12 \cdot AG \cdot BC = \frac12 (2 - \sqrt{3})(\sqrt{3} - 1)$
$[\text{12-sided outer polygon}] = [\text{cross}] + 8 \cdot [ABC] = 4\sqrt{3} - 3 + 8 \cdot \frac12 (2 - \sqrt{3})(\sqrt{3} - 1) = 4\sqrt{3} - 3 + 4 (3\sqrt{3} - 5) = 16\sqrt{3}-23$

 \[16+3-23=\boxed{\textbf{(B)}\ -4}\]




\textbf{Solution 4}

We calculate the area as the area of the red octagon minus the four purple congruent triangles:
                    \begin{figure}[H]
                        \centering
                        \includegraphics[width=0.5\textwidth]{Solution/AMC_12_12B_2022_25_Solution_Figure_5.png}
                    \end{figure}
                    We first find the important angles in the figure. We note that 2 adjacent hexagons are rotated $90^\circ$ with respect to the other, so the angles between any sides is $150^\circ$. In particular, as the purple triangles are isosceles, they have angles $150^\circ,15^\circ$, and $15^\circ$, and the octagon is equiangular (all its angles are $135^\circ$). Thus, we can draw a square around the octagon, and we note that the ``cut out" triangles are all isosceles right triangles.

 Now, we calculate the side length of the square. Note that the hexagon has a height of $\sqrt 3$, so the length of a side of the square is $2\sqrt 3-1$. In particular, the horizontal/vertical sides of the octagon have length $1$, so the legs of the isosceles triangles are
 \[\frac{2\sqrt3-1-1}2=\sqrt3-1\]Thus, the area of the octagon is
 \[(2\sqrt3-1)^2-4\cdot\frac 12(\sqrt3-1)^2=5\]Now, we calculate the area of one of the four isosceles triangles. The base of the triangle is $(\sqrt 3-1)\sqrt 2$, so the area is
 \[\frac 14\mathrm{base length}^2\cdot\tan(\mathrm{base angle})=\frac 14((\sqrt3-1)\sqrt2)^2\cdot\tan15^\circ=\frac 14(8-4\sqrt3)(2-\sqrt3)=7-4\sqrt 3\]Thus, the area of the dodecagon is
 \[5-4(7-4\sqrt3)=16\sqrt3-23\]Thus the answer is $16+3-23=-4$, or $\boxed{\textbf{(B)}}$.

 



\textbf{Solution 5 (No Trig, No Coordbash)}


                    \begin{figure}[H]
                        \centering
                        \includegraphics[width=0.5\textwidth]{Solution/AMC_12_12B_2022_25_Solution_Figure_6.jpeg}
                    \end{figure}
                    Note that each of the green sections is a rectangle, so its interior angles are all $90^{\circ}.$ Since $\angle{ABC}=120^{\circ}$, every one of the orange sections is a $30-60-90$ right triangle. 

 Define $x$ to be the distance from the corner of the square with side length $1$ to the corner of the larger blue square. Due to the sides of the two squares being parallel to each other, the large blue triangle is a $45-45-90$ right triangle. By $AAA$ similarity, the smaller blue triangles are also $45-45-90,$ and have side lengths of $\frac{x\sqrt{2}}{2}, \frac{x\sqrt{2}}{2},$ and $x$. By $30-60-90$ triangle relations, the largest altitude of the orange triangle is $\frac{x\sqrt{6}}{2}.$ 

 Now, we can find the height of the hexagon to obtain an equation in terms of $x$. Consider a hexagon with side length $1$, where point $P$ is the foot of the perpendicular dropped from $B$, bisecting $\angle{ABC}$:
                    \begin{figure}[H]
                        \centering
                        \includegraphics[width=0.5\textwidth]{Solution/AMC_12_12B_2022_25_Solution_Figure_7.png}
                    \end{figure}
                    Note that triangles $\triangle{APB}$ and $\triangle{CPB}$ are congruent $30-60-90$ triangles, by SAS congruence. Since the side length of this hexagon is $1$, the length of $AP$ is $\frac{1}{2}\cdot \sqrt{3} =\frac{\sqrt{3}}{2}$, by $30-60-90$ triangle relations. The height of the hexagon is twice this value, or $\sqrt{3}.$ 

 The height is also equal to the sum of the values along the long blue line, in the first diagram. Therefore, $1+\frac{x\sqrt{2}}{2}+\frac{x\sqrt{6}}{2}=\sqrt{3}.$ Solving and rationalizing, $x=2\sqrt{2}-\sqrt{6}.$

 The area of the dodecagon is equal to the sum of the areas of the four rectangles, eight orange triangles, and purple square. In terms of $x$, this is 
 \[8\cdot \frac{1}{2}\cdot \frac{x\sqrt{2}}{2}\cdot \frac{x\sqrt{6}}{{2}}+4\cdot 1\cdot\frac{x\sqrt{6}}{2}+(x\sqrt{2}+1)^2.\]

 Plugging in $x=2\sqrt{2}-\sqrt{6}$, the area of the dodecagon is $16\sqrt{3}-23$. Therefore, the answer is $16+3-23=$ $\boxed{\textbf{(B) }-4}.$

 -Benedict T (countmath1)

 


 



\textbf{Solution 6 (PIE bash)}

Obviously, each of the hexagon have area of $\frac{3\sqrt{3}}{2}$, which gives a total of $6\sqrt{3}$ (without accounting for the overcounted overlapped areas). Now we calculate the area of the intersection of two hexagons. There are two cases : 

 
\textbf{Case 1}: The hexagons are on the opposite sides

 In this case, we see that the two hexagons' intersection is the middle square plus two triangles on its side. The square have area $1$. Each of the triangles are $120-30-30$, and hence its short side is $\frac{\sqrt{3}}{3}$ because its longer side is $1$. Therefore, each of the small triangles have area of $\frac{\sqrt{3}}{12}$, so each of the intersections have $1+\frac{\sqrt{3}}{6}$ area. Since there are $2$ pairs of hexagons on opposite side of each other, that gives a total area of $2+\frac{\sqrt{3}}{3}$.

 
\textbf{Case 2}: The hexagons are adjacent
                    \begin{figure}[H]
                        \centering
                        \includegraphics[width=0.5\textwidth]{Solution/AMC_12_12B_2022_25_Solution_Figure_8.png}
                    \end{figure}
                    

 In this case, their intersection is a convex hexagon. We cut it in half and we're left with a quadrilateral with two sides with length one (because it shares these sides with one of the hexagons, and for this reason the angle between the two sides is $120$ degrees) and one short side and one long side. Here, we divide the area even further. We draw a line that divide the quadrilateral into a triangle with side lengths $1$, $1$, and $\sqrt{3}$(because it is $30-30-120$) and another triangle with unknown side lengths. Simple angle chasing yields the triangle with unknown side lengths is a $15-75-90$ triangle. If you remembered your sines and cosines, you know that $\sin{15}=\frac{\sqrt{6}-\sqrt{2}}{4}$ or you can calculate it by using the fact that $\sin{15}=\sin{45-30}$. Also, $\sin{75}=\frac{\sqrt{6}+\sqrt{2}}{4}$. Now, because the $15-75-90$ triangle have one side with length $\sqrt{3}$(it shares the side with the other triangle) and it corresponds to the 75 degrees angle, we have 

 
 \[\frac{\sqrt{3}}{\frac{\sqrt{6}+\sqrt{2}}{4}}=\frac{x}{\frac{\sqrt{6}-\sqrt{2}}{4}}\]

 \[\implies x=\frac{\sqrt{3}{(\sqrt{6}-\sqrt{2})}}{\sqrt{6}+\sqrt{2}}=2\sqrt{3}-3\]

 where $x$ denote the shorter leg of the $15-75-90$ triangle. Hence, that triangle have area 

 
 \[\frac{\sqrt{3}}{2} x=\frac{\sqrt{3}}{2}(2\sqrt{3}-3)=\frac{6-3\sqrt{3}}{2}\]. 

 The isosceles $120-30-30$ triangle obviously have area $\frac{\sqrt{3}}{4}$, and that gives a total of $\frac{12-5\sqrt{3}}{4}$. We multiply that by $2$ because $\frac{12-5\sqrt{3}}{4}$ is only half of the intersection of the hexagons. We multiply the result by $4$ because there are $4$ pairs of adjacent hexagons. That gives $24-10\sqrt{3}$, which is it for Case 2. 

 
Now that we got the total area of the intersection of two hexagons, we find the area of the intersection of three hexagons. Note that the intersection of any three hexagons is just the middle square plus one small triangle on the side. Again, the middle square have area $1$ and the triangle have area $\frac{\sqrt{3}}{12}$, so the total is $1+\frac{\sqrt{3}}{12}$. We multiply that by $4\choose3$ to get $4+\frac{\sqrt{3}}{3}$.

 
Then, we calculate the area of the intersection of all four hexagons, which is just the middle square with area $1$. 

 
Finally, we apply PIE. The total area is 

 
 \[6\sqrt{3}-[(2+\frac{\sqrt{3}}{3})+(24-10\sqrt{3})]+(4+\frac{\sqrt{3}}{3})-1=16\sqrt{3}-23\],

 so our final answer is $16+3-23=\boxed{(B)-4}$. $\blacksquare$ ~\href{https://artofproblemsolving.com/wiki/index.php?title=Ddk001}{Ddk001}

 


\end{enumerate}